% Paper v2 (MVP1 + MVP2): v1 plus Q1a topic, Q1b deadline horizon, Q1c volume. Run 142083.
% Every measured number is a macro from numbers.tex (generated, do not edit).
%
% Source key for the "% src:" comments (this key is the Sources table of this file):
%   S-PAPER-1 proposal_md/polymarket_proposal_v3.md
%   S-PAPER-2 tips/polymarket_plan_detailed.md
%   S-PAPER-3 analysis/prereg_mvp1.md
%   S-PAPER-4 analysis/prereg_mvp2.md
%   S-PAPER-5 runs/142083_mvp2_full_v2r3 (reports, verdicts.json) via numbers.tex
%   S-PAPER-6 analysis/review_mvp1_mvp2.md
%   S-PAPER-7 tips/polymarket_class_questions_tips.html
%   S-PAPER-8 PLAN.md, RQ.md, DECISIONS.md, STATUS.md
%   S-PAPER-9 analysis/prereg_robustness.md, runs/142123_robust_full_v4 (reports r1, r2) via numbers.tex
\documentclass[11pt]{article}

\usepackage[margin=0.9in]{geometry}
\usepackage[T1]{fontenc}
\usepackage[utf8]{inputenc}
\usepackage{lmodern}
\usepackage{microtype}
\usepackage{amsmath}
\usepackage{booktabs}
\usepackage{tabularx}
\usepackage[numbers]{natbib}
\usepackage{figkit}
\usepackage[hidelinks]{hyperref}
\usepackage[capitalise,noabbrev]{cleveref}

\captionsetup{font=small,labelfont=bf}
\setlength{\emergencystretch}{2em}
\setlength{\intextsep}{6pt plus 2pt minus 2pt}
\setlength{\textfloatsep}{8pt plus 2pt minus 2pt}
\setlength{\floatsep}{6pt plus 2pt minus 2pt}
\setlength{\abovecaptionskip}{4pt}

% Visible placeholder for text or numbers the sources do not give.
\DeclareRobustCommand{\todo}[1]{\textcolor{red}{\textbf{[TODO: #1]}}}

% Generated by code/make_paper_numbers.py from runs/142083_mvp2_full_v2r3. Do not edit by hand.
% run directory name
\newcommand{\RunId}{142083\_mvp2\_full\_v2r3}
% manifest_q1_calibration.json: code_version
\newcommand{\CodeVersion}{v3}
% manifest_q1_calibration.json: code_sha (hash of the script q1_calibration.py alone, first 12 hex, not of the code set)
\newcommand{\CodeSha}{8efcddf9ef94}
% code.sha256: sha256 of the file that lists the hash of every script (first 12 hex)
\newcommand{\CodeSetSha}{950ec496595b}
% path of the per-script hash list
\newcommand{\CodeShaFile}{runs/142083\_mvp2\_full\_v2r3/code.sha256}
% manifest_q1_calibration.json: env.slurm_job_id
\newcommand{\SlurmJobId}{142083}
% verdicts.json: comparisons.Q1.verdict
\newcommand{\VerdictQOne}{Prices are not calibrated as a whole at this n.}
% verdicts.json: comparisons.Q1a.verdict
\newcommand{\VerdictQOneA}{Accuracy depends on the topic at this n.}
% verdicts.json: comparisons.Q1b.verdict
\newcommand{\VerdictQOneB}{Better near the deadline.}
% verdicts.json: comparisons.Q1c.verdict
\newcommand{\VerdictQOneC}{No evidence that busier markets are more accurate at this n.}
% verdicts.json: alpha
\newcommand{\Alpha}{\ensuremath{0.05}}
% tables_summary.json: n_markets_in_scope
\newcommand{\NMarketsInScope}{\ensuremath{688{,}131}}
% tables_summary.json: n_kept
\newcommand{\NMarketsKept}{\ensuremath{47{,}948}}
% tables_summary.json: n_kept / n_markets_in_scope
\newcommand{\PctMarketsKept}{\ensuremath{7.0}\%}
% tables_summary.json: n_kept_no_min_filter
\newcommand{\NMarketsKeptNoMin}{\ensuremath{74{,}474}}
% tables_summary.json: drop_reasons.below_min_trades
\newcommand{\NBelowMinTrades}{\ensuremath{26{,}526}}
% tables_summary.json: drop_reasons.below_min_trades / n_markets_in_scope
\newcommand{\PctBelowMinTrades}{\ensuremath{3.9}\%}
% tables_summary.json: drop_reasons.stale_price
\newcommand{\NStalePrice}{\ensuremath{80{,}230}}
% tables_summary.json: drop_reasons.stale_price / n_markets_in_scope
\newcommand{\PctStalePrice}{\ensuremath{11.7}\%}
% tables_summary.json: drop_reasons.updown
\newcommand{\NUpdown}{\ensuremath{255{,}420}}
% tables_summary.json: drop_reasons.updown / n_markets_in_scope
\newcommand{\PctUpdown}{\ensuremath{37.1}\%}
% tables_summary.json: drop_reasons.kept
\newcommand{\NKept}{\ensuremath{47{,}948}}
% tables_summary.json: drop_reasons.kept / n_markets_in_scope
\newcommand{\PctKept}{\ensuremath{7.0}\%}
% tables_summary.json: drop_reasons.no_outcome
\newcommand{\NNoOutcome}{\ensuremath{62{,}718}}
% tables_summary.json: drop_reasons.no_outcome / n_markets_in_scope
\newcommand{\PctNoOutcome}{\ensuremath{9.1}\%}
% tables_summary.json: drop_reasons.no_price
\newcommand{\NNoPrice}{\ensuremath{215{,}289}}
% tables_summary.json: drop_reasons.no_price / n_markets_in_scope
\newcommand{\PctNoPrice}{\ensuremath{31.3}\%}
% tables_summary.json: n_trades_in_kept_markets
\newcommand{\NTradesKeptMarkets}{\ensuremath{63{,}023{,}987}}
% tables_summary.json: n_wallet_role_rows
\newcommand{\NWalletRoleRows}{\ensuremath{2{,}117{,}412}}
% tables_summary.json: trade_size_sample_rows
\newcommand{\NTradeSizeSample}{\ensuremath{1{,}000{,}000}}
% tables_summary.json: n_kept_open_ge_7d
\newcommand{\NKeptOpenSevenDays}{\ensuremath{21{,}816}}
% tables_summary.json: n_kept_with_p_7d
\newcommand{\NKeptWithPriceSevenDays}{\ensuremath{14{,}169}}
% q1_results.json: main_with_min_filter.n_markets
\newcommand{\QOneN}{\ensuremath{47{,}948}}
% q1_results.json: main_with_min_filter.chi2_gof.chi2
\newcommand{\QOneChiSq}{\ensuremath{68.0}}
% q1_results.json: main_with_min_filter.chi2_gof.df
\newcommand{\QOneDf}{\ensuremath{10}}
% q1_results.json: main_with_min_filter.chi2_gof.p
\newcommand{\QOneP}{\ensuremath{1.1\times10^{-10}}}
% q1_results.json: main_with_min_filter.chi2_gof.effect_w
\newcommand{\QOneW}{\ensuremath{0.038}}
% q1_results.json: main_with_min_filter.brier
\newcommand{\QOneBrier}{\ensuremath{0.107}}
% q1_results.json: main_with_min_filter.log_loss
\newcommand{\QOneLogLoss}{\ensuremath{0.323}}
% q1_results.json: main_with_min_filter.mean_entropy_bits
\newcommand{\QOneEntropy}{\ensuremath{0.475}}
% q1_results.json: main_with_min_filter.class_balance_yes
\newcommand{\QOneYesRate}{\ensuremath{36.6}\%}
% q1_results.json: main_with_min_filter.confusion_at_0.5.accuracy
\newcommand{\QOneAccuracy}{\ensuremath{83.4}\%}
% q1_results.json: main_with_min_filter.confusion_at_0.5.precision
\newcommand{\QOnePrecision}{\ensuremath{77.7}\%}
% q1_results.json: main_with_min_filter.confusion_at_0.5.recall
\newcommand{\QOneRecall}{\ensuremath{76.8}\%}
% q1_results.json: main_with_min_filter.ks_yes_vs_no.D
\newcommand{\QOneKsD}{\ensuremath{0.673}}
% q1_results.json: main_with_min_filter.bins_reject_bh
\newcommand{\QOneBinsBh}{\ensuremath{5}}
% q1_results.json: main_with_min_filter.bins_reject_bonferroni
\newcommand{\QOneBinsBonferroni}{\ensuremath{4}}
% q1_results.json: no_min_filter.n_markets
\newcommand{\QOneNoMinN}{\ensuremath{74{,}474}}
% q1_results.json: no_min_filter.brier
\newcommand{\QOneNoMinBrier}{\ensuremath{0.134}}
% q1_results.json: no_min_filter.chi2_gof.chi2
\newcommand{\QOneNoMinChiSq}{\ensuremath{214.2}}
% q1_results.json: no_min_filter.chi2_gof.p
\newcommand{\QOneNoMinP}{\ensuremath{1.7\times10^{-40}}}
% q1_results.json: no_min_filter.chi2_gof.effect_w
\newcommand{\QOneNoMinW}{\ensuremath{0.054}}
% q1_results.json: min_trades
\newcommand{\MinTrades}{\ensuremath{10}}
% q1_bins.csv: bin=0 n
\newcommand{\BinZeroTenN}{\ensuremath{18{,}209}}
% q1_bins.csv: bin=0 mean_price
\newcommand{\BinZeroTenPrice}{\ensuremath{1.8}\%}
% q1_bins.csv: bin=0 observed_rate
\newcommand{\BinZeroTenObs}{\ensuremath{1.4}\%}
% q1_bins.csv: bin=0 diff_obs_minus_price
\newcommand{\BinZeroTenDiff}{\ensuremath{-0.46}\%}
% q1_bins.csv: bin=0 ci_low
\newcommand{\BinZeroTenCiLow}{\ensuremath{1.2}\%}
% q1_bins.csv: bin=0 ci_high
\newcommand{\BinZeroTenCiHigh}{\ensuremath{1.5}\%}
% q1_bins.csv: bin=0 binom_p
\newcommand{\BinZeroTenP}{\ensuremath{1.6\times10^{-6}}}
% q1_bins.csv: bin=0 p_bh
\newcommand{\BinZeroTenBhP}{\ensuremath{1.6\times10^{-5}}}
% q1_bins.csv: bin=1 n
\newcommand{\BinTenTwentyN}{\ensuremath{2{,}703}}
% q1_bins.csv: bin=1 mean_price
\newcommand{\BinTenTwentyPrice}{\ensuremath{14.8}\%}
% q1_bins.csv: bin=1 observed_rate
\newcommand{\BinTenTwentyObs}{\ensuremath{12.5}\%}
% q1_bins.csv: bin=1 diff_obs_minus_price
\newcommand{\BinTenTwentyDiff}{\ensuremath{-2.30}\%}
% q1_bins.csv: bin=1 ci_low
\newcommand{\BinTenTwentyCiLow}{\ensuremath{11.3}\%}
% q1_bins.csv: bin=1 ci_high
\newcommand{\BinTenTwentyCiHigh}{\ensuremath{13.8}\%}
% q1_bins.csv: bin=1 binom_p
\newcommand{\BinTenTwentyP}{\ensuremath{6.5\times10^{-4}}}
% q1_bins.csv: bin=1 p_bh
\newcommand{\BinTenTwentyBhP}{\ensuremath{0.002}}
% q1_bins.csv: bin=2 n
\newcommand{\BinTwentyThirtyN}{\ensuremath{2{,}543}}
% q1_bins.csv: bin=2 mean_price
\newcommand{\BinTwentyThirtyPrice}{\ensuremath{24.9}\%}
% q1_bins.csv: bin=2 observed_rate
\newcommand{\BinTwentyThirtyObs}{\ensuremath{23.9}\%}
% q1_bins.csv: bin=2 diff_obs_minus_price
\newcommand{\BinTwentyThirtyDiff}{\ensuremath{-1.00}\%}
% q1_bins.csv: bin=2 ci_low
\newcommand{\BinTwentyThirtyCiLow}{\ensuremath{22.3}\%}
% q1_bins.csv: bin=2 ci_high
\newcommand{\BinTwentyThirtyCiHigh}{\ensuremath{25.6}\%}
% q1_bins.csv: bin=2 binom_p
\newcommand{\BinTwentyThirtyP}{\ensuremath{0.252}}
% q1_bins.csv: bin=2 p_bh
\newcommand{\BinTwentyThirtyBhP}{\ensuremath{0.315}}
% q1_bins.csv: bin=3 n
\newcommand{\BinThirtyFortyN}{\ensuremath{2{,}769}}
% q1_bins.csv: bin=3 mean_price
\newcommand{\BinThirtyFortyPrice}{\ensuremath{34.9}\%}
% q1_bins.csv: bin=3 observed_rate
\newcommand{\BinThirtyFortyObs}{\ensuremath{33.2}\%}
% q1_bins.csv: bin=3 diff_obs_minus_price
\newcommand{\BinThirtyFortyDiff}{\ensuremath{-1.72}\%}
% q1_bins.csv: bin=3 ci_low
\newcommand{\BinThirtyFortyCiLow}{\ensuremath{31.5}\%}
% q1_bins.csv: bin=3 ci_high
\newcommand{\BinThirtyFortyCiHigh}{\ensuremath{35.0}\%}
% q1_bins.csv: bin=3 binom_p
\newcommand{\BinThirtyFortyP}{\ensuremath{0.058}}
% q1_bins.csv: bin=3 p_bh
\newcommand{\BinThirtyFortyBhP}{\ensuremath{0.097}}
% q1_bins.csv: bin=4 n
\newcommand{\BinFortyFiftyN}{\ensuremath{3{,}859}}
% q1_bins.csv: bin=4 mean_price
\newcommand{\BinFortyFiftyPrice}{\ensuremath{45.0}\%}
% q1_bins.csv: bin=4 observed_rate
\newcommand{\BinFortyFiftyObs}{\ensuremath{43.8}\%}
% q1_bins.csv: bin=4 diff_obs_minus_price
\newcommand{\BinFortyFiftyDiff}{\ensuremath{-1.23}\%}
% q1_bins.csv: bin=4 ci_low
\newcommand{\BinFortyFiftyCiLow}{\ensuremath{42.3}\%}
% q1_bins.csv: bin=4 ci_high
\newcommand{\BinFortyFiftyCiHigh}{\ensuremath{45.4}\%}
% q1_bins.csv: bin=4 binom_p
\newcommand{\BinFortyFiftyP}{\ensuremath{0.128}}
% q1_bins.csv: bin=4 p_bh
\newcommand{\BinFortyFiftyBhP}{\ensuremath{0.183}}
% q1_bins.csv: bin=5 n
\newcommand{\BinFiftySixtyN}{\ensuremath{4{,}240}}
% q1_bins.csv: bin=5 mean_price
\newcommand{\BinFiftySixtyPrice}{\ensuremath{54.1}\%}
% q1_bins.csv: bin=5 observed_rate
\newcommand{\BinFiftySixtyObs}{\ensuremath{53.5}\%}
% q1_bins.csv: bin=5 diff_obs_minus_price
\newcommand{\BinFiftySixtyDiff}{\ensuremath{-0.58}\%}
% q1_bins.csv: bin=5 ci_low
\newcommand{\BinFiftySixtyCiLow}{\ensuremath{52.0}\%}
% q1_bins.csv: bin=5 ci_high
\newcommand{\BinFiftySixtyCiHigh}{\ensuremath{55.0}\%}
% q1_bins.csv: bin=5 binom_p
\newcommand{\BinFiftySixtyP}{\ensuremath{0.450}}
% q1_bins.csv: bin=5 p_bh
\newcommand{\BinFiftySixtyBhP}{\ensuremath{0.500}}
% q1_bins.csv: bin=6 n
\newcommand{\BinSixtySeventyN}{\ensuremath{2{,}553}}
% q1_bins.csv: bin=6 mean_price
\newcommand{\BinSixtySeventyPrice}{\ensuremath{64.4}\%}
% q1_bins.csv: bin=6 observed_rate
\newcommand{\BinSixtySeventyObs}{\ensuremath{61.7}\%}
% q1_bins.csv: bin=6 diff_obs_minus_price
\newcommand{\BinSixtySeventyDiff}{\ensuremath{-2.79}\%}
% q1_bins.csv: bin=6 ci_low
\newcommand{\BinSixtySeventyCiLow}{\ensuremath{59.8}\%}
% q1_bins.csv: bin=6 ci_high
\newcommand{\BinSixtySeventyCiHigh}{\ensuremath{63.5}\%}
% q1_bins.csv: bin=6 binom_p
\newcommand{\BinSixtySeventyP}{\ensuremath{0.004}}
% q1_bins.csv: bin=6 p_bh
\newcommand{\BinSixtySeventyBhP}{\ensuremath{0.009}}
% q1_bins.csv: bin=7 n
\newcommand{\BinSeventyEightyN}{\ensuremath{2{,}301}}
% q1_bins.csv: bin=7 mean_price
\newcommand{\BinSeventyEightyPrice}{\ensuremath{74.6}\%}
% q1_bins.csv: bin=7 observed_rate
\newcommand{\BinSeventyEightyObs}{\ensuremath{71.3}\%}
% q1_bins.csv: bin=7 diff_obs_minus_price
\newcommand{\BinSeventyEightyDiff}{\ensuremath{-3.29}\%}
% q1_bins.csv: bin=7 ci_low
\newcommand{\BinSeventyEightyCiLow}{\ensuremath{69.4}\%}
% q1_bins.csv: bin=7 ci_high
\newcommand{\BinSeventyEightyCiHigh}{\ensuremath{73.1}\%}
% q1_bins.csv: bin=7 binom_p
\newcommand{\BinSeventyEightyP}{\ensuremath{3.6\times10^{-4}}}
% q1_bins.csv: bin=7 p_bh
\newcommand{\BinSeventyEightyBhP}{\ensuremath{0.002}}
% q1_bins.csv: bin=8 n
\newcommand{\BinEightyNinetyN}{\ensuremath{2{,}133}}
% q1_bins.csv: bin=8 mean_price
\newcommand{\BinEightyNinetyPrice}{\ensuremath{84.5}\%}
% q1_bins.csv: bin=8 observed_rate
\newcommand{\BinEightyNinetyObs}{\ensuremath{82.7}\%}
% q1_bins.csv: bin=8 diff_obs_minus_price
\newcommand{\BinEightyNinetyDiff}{\ensuremath{-1.88}\%}
% q1_bins.csv: bin=8 ci_low
\newcommand{\BinEightyNinetyCiLow}{\ensuremath{81.0}\%}
% q1_bins.csv: bin=8 ci_high
\newcommand{\BinEightyNinetyCiHigh}{\ensuremath{84.2}\%}
% q1_bins.csv: bin=8 binom_p
\newcommand{\BinEightyNinetyP}{\ensuremath{0.018}}
% q1_bins.csv: bin=8 p_bh
\newcommand{\BinEightyNinetyBhP}{\ensuremath{0.036}}
% q1_bins.csv: bin=9 n
\newcommand{\BinNinetyHundredN}{\ensuremath{6{,}638}}
% q1_bins.csv: bin=9 mean_price
\newcommand{\BinNinetyHundredPrice}{\ensuremath{97.7}\%}
% q1_bins.csv: bin=9 observed_rate
\newcommand{\BinNinetyHundredObs}{\ensuremath{97.8}\%}
% q1_bins.csv: bin=9 diff_obs_minus_price
\newcommand{\BinNinetyHundredDiff}{\ensuremath{0.03}\%}
% q1_bins.csv: bin=9 ci_low
\newcommand{\BinNinetyHundredCiLow}{\ensuremath{97.4}\%}
% q1_bins.csv: bin=9 ci_high
\newcommand{\BinNinetyHundredCiHigh}{\ensuremath{98.1}\%}
% q1_bins.csv: bin=9 binom_p
\newcommand{\BinNinetyHundredP}{\ensuremath{0.902}}
% q1_bins.csv: bin=9 p_bh
\newcommand{\BinNinetyHundredBhP}{\ensuremath{0.902}}
% q1a_results.json: chi_square_independence.n
\newcommand{\QOneAN}{\ensuremath{47{,}948}}
% q1a_results.json: chi_square_independence.chi2
\newcommand{\QOneAChiSq}{\ensuremath{3203.8}}
% q1a_results.json: chi_square_independence.df
\newcommand{\QOneADf}{\ensuremath{7}}
% q1a_results.json: chi_square_independence.p
\newcommand{\QOneAP}{\ensuremath{<10^{-300}}}
% q1a_results.json: chi_square_independence.cramers_v
\newcommand{\QOneAV}{\ensuremath{0.258}}
% q1a_results.json: chi_square_independence.min_expected_count
\newcommand{\QOneAMinExpected}{\ensuremath{109.8}}
% q1a_results.json: chi_square_independence.n_categories
\newcommand{\QOneANCategories}{\ensuremath{8}}
% q1a_results.json: bin_tests.n_tested
\newcommand{\QOneABinsTested}{\ensuremath{80}}
% q1a_results.json: bin_tests.reject_bh
\newcommand{\QOneABinsBh}{\ensuremath{9}}
% q1a_results.json: bin_tests.reject_bonferroni
\newcommand{\QOneABinsBonferroni}{\ensuremath{2}}
% q1a_results.json: brier_tests.anova_F
\newcommand{\QOneAAnovaF}{\ensuremath{837.7}}
% q1a_results.json: brier_tests.anova_p
\newcommand{\QOneAAnovaP}{\ensuremath{<10^{-300}}}
% q1a_results.json: brier_tests.eta_squared
\newcommand{\QOneAEtaSq}{\ensuremath{0.109}}
% q1a_results.json: brier_tests.perm_p
\newcommand{\QOneAPermP}{\ensuremath{1.0\times10^{-3}}}
% q1a_results.json: brier_tests.n_perm
\newcommand{\QOneANPerm}{\ensuremath{1{,}000}}
% q1_results.json: main_with_min_filter.confusion_at_0.5.accuracy (share of kept markets whose p>0.5 side won)
\newcommand{\CatAllFavWin}{\ensuremath{83.4}\%}
% q1a_contingency.csv: category_group=Sports n
\newcommand{\CatSportsN}{\ensuremath{19{,}162}}
% q1a_contingency.csv: category_group=Sports favorite_win_rate
\newcommand{\CatSportsFavWin}{\ensuremath{72.2}\%}
% q1a_contingency.csv: category_group=Sports mean_brier
\newcommand{\CatSportsBrier}{\ensuremath{0.174}}
% q1a_contingency.csv: category_group=Crypto n
\newcommand{\CatCryptoN}{\ensuremath{11{,}888}}
% q1a_contingency.csv: category_group=Crypto favorite_win_rate
\newcommand{\CatCryptoFavWin}{\ensuremath{94.8}\%}
% q1a_contingency.csv: category_group=Crypto mean_brier
\newcommand{\CatCryptoBrier}{\ensuremath{0.037}}
% q1a_contingency.csv: category_group=Politics n
\newcommand{\CatPoliticsN}{\ensuremath{6{,}922}}
% q1a_contingency.csv: category_group=Politics favorite_win_rate
\newcommand{\CatPoliticsFavWin}{\ensuremath{86.6}\%}
% q1a_contingency.csv: category_group=Politics mean_brier
\newcommand{\CatPoliticsBrier}{\ensuremath{0.092}}
% q1a_contingency.csv: category_group=Finance n
\newcommand{\CatFinanceN}{\ensuremath{5{,}146}}
% q1a_contingency.csv: category_group=Finance favorite_win_rate
\newcommand{\CatFinanceFavWin}{\ensuremath{90.6}\%}
% q1a_contingency.csv: category_group=Finance mean_brier
\newcommand{\CatFinanceBrier}{\ensuremath{0.067}}
% q1a_contingency.csv: category_group=Other n
\newcommand{\CatOtherN}{\ensuremath{2{,}274}}
% q1a_contingency.csv: category_group=Other favorite_win_rate
\newcommand{\CatOtherFavWin}{\ensuremath{85.6}\%}
% q1a_contingency.csv: category_group=Other mean_brier
\newcommand{\CatOtherBrier}{\ensuremath{0.096}}
% q1a_contingency.csv: category_group=Culture n
\newcommand{\CatCultureN}{\ensuremath{1{,}074}}
% q1a_contingency.csv: category_group=Culture favorite_win_rate
\newcommand{\CatCultureFavWin}{\ensuremath{90.4}\%}
% q1a_contingency.csv: category_group=Culture mean_brier
\newcommand{\CatCultureBrier}{\ensuremath{0.066}}
% q1a_contingency.csv: category_group=Sci-Tech n
\newcommand{\CatSciTechN}{\ensuremath{819}}
% q1a_contingency.csv: category_group=Sci-Tech favorite_win_rate
\newcommand{\CatSciTechFavWin}{\ensuremath{88.6}\%}
% q1a_contingency.csv: category_group=Sci-Tech mean_brier
\newcommand{\CatSciTechBrier}{\ensuremath{0.079}}
% q1a_contingency.csv: category_group=Price Action n
\newcommand{\CatPriceActionN}{\ensuremath{663}}
% q1a_contingency.csv: category_group=Price Action favorite_win_rate
\newcommand{\CatPriceActionFavWin}{\ensuremath{92.9}\%}
% q1a_contingency.csv: category_group=Price Action mean_brier
\newcommand{\CatPriceActionBrier}{\ensuremath{0.049}}
% q1a_bins_by_category.csv: category=Crypto bin=0 n
\newcommand{\CatCryptoBinZeroTenN}{\ensuremath{6{,}546}}
% q1a_bins_by_category.csv: category=Crypto bin=0 mean_price
\newcommand{\CatCryptoBinZeroTenPrice}{\ensuremath{1.1}\%}
% q1a_bins_by_category.csv: category=Crypto bin=0 observed_rate
\newcommand{\CatCryptoBinZeroTenObs}{\ensuremath{0.7}\%}
% q1a_bins_by_category.csv: category=Crypto bin=0 diff_obs_minus_price
\newcommand{\CatCryptoBinZeroTenDiff}{\ensuremath{-0.44}\%}
% q1a_bins_by_category.csv: category=Crypto bin=0 p_bh
\newcommand{\CatCryptoBinZeroTenBhP}{\ensuremath{0.017}}
% q1a_bins_by_category.csv: category=Crypto bin=1 n
\newcommand{\CatCryptoBinTenTwentyN}{\ensuremath{399}}
% q1a_bins_by_category.csv: category=Crypto bin=1 mean_price
\newcommand{\CatCryptoBinTenTwentyPrice}{\ensuremath{14.5}\%}
% q1a_bins_by_category.csv: category=Crypto bin=1 observed_rate
\newcommand{\CatCryptoBinTenTwentyObs}{\ensuremath{8.8}\%}
% q1a_bins_by_category.csv: category=Crypto bin=1 diff_obs_minus_price
\newcommand{\CatCryptoBinTenTwentyDiff}{\ensuremath{-5.75}\%}
% q1a_bins_by_category.csv: category=Crypto bin=1 p_bh
\newcommand{\CatCryptoBinTenTwentyBhP}{\ensuremath{0.021}}
% q1a_bins_by_category.csv: category=Culture bin=0 n
\newcommand{\CatCultureBinZeroTenN}{\ensuremath{552}}
% q1a_bins_by_category.csv: category=Culture bin=0 mean_price
\newcommand{\CatCultureBinZeroTenPrice}{\ensuremath{1.8}\%}
% q1a_bins_by_category.csv: category=Culture bin=0 observed_rate
\newcommand{\CatCultureBinZeroTenObs}{\ensuremath{1.4}\%}
% q1a_bins_by_category.csv: category=Culture bin=0 diff_obs_minus_price
\newcommand{\CatCultureBinZeroTenDiff}{\ensuremath{-0.35}\%}
% q1a_bins_by_category.csv: category=Culture bin=0 p_bh
\newcommand{\CatCultureBinZeroTenBhP}{\ensuremath{0.979}}
% q1a_bins_by_category.csv: category=Culture bin=1 n
\newcommand{\CatCultureBinTenTwentyN}{\ensuremath{61}}
% q1a_bins_by_category.csv: category=Culture bin=1 mean_price
\newcommand{\CatCultureBinTenTwentyPrice}{\ensuremath{14.8}\%}
% q1a_bins_by_category.csv: category=Culture bin=1 observed_rate
\newcommand{\CatCultureBinTenTwentyObs}{\ensuremath{8.2}\%}
% q1a_bins_by_category.csv: category=Culture bin=1 diff_obs_minus_price
\newcommand{\CatCultureBinTenTwentyDiff}{\ensuremath{-6.60}\%}
% q1a_bins_by_category.csv: category=Culture bin=1 p_bh
\newcommand{\CatCultureBinTenTwentyBhP}{\ensuremath{0.548}}
% q1a_bins_by_category.csv: category=Finance bin=0 n
\newcommand{\CatFinanceBinZeroTenN}{\ensuremath{2{,}291}}
% q1a_bins_by_category.csv: category=Finance bin=0 mean_price
\newcommand{\CatFinanceBinZeroTenPrice}{\ensuremath{1.4}\%}
% q1a_bins_by_category.csv: category=Finance bin=0 observed_rate
\newcommand{\CatFinanceBinZeroTenObs}{\ensuremath{0.8}\%}
% q1a_bins_by_category.csv: category=Finance bin=0 diff_obs_minus_price
\newcommand{\CatFinanceBinZeroTenDiff}{\ensuremath{-0.62}\%}
% q1a_bins_by_category.csv: category=Finance bin=0 p_bh
\newcommand{\CatFinanceBinZeroTenBhP}{\ensuremath{0.079}}
% q1a_bins_by_category.csv: category=Finance bin=1 n
\newcommand{\CatFinanceBinTenTwentyN}{\ensuremath{164}}
% q1a_bins_by_category.csv: category=Finance bin=1 mean_price
\newcommand{\CatFinanceBinTenTwentyPrice}{\ensuremath{14.5}\%}
% q1a_bins_by_category.csv: category=Finance bin=1 observed_rate
\newcommand{\CatFinanceBinTenTwentyObs}{\ensuremath{14.0}\%}
% q1a_bins_by_category.csv: category=Finance bin=1 diff_obs_minus_price
\newcommand{\CatFinanceBinTenTwentyDiff}{\ensuremath{-0.43}\%}
% q1a_bins_by_category.csv: category=Finance bin=1 p_bh
\newcommand{\CatFinanceBinTenTwentyBhP}{\ensuremath{1.000}}
% q1a_bins_by_category.csv: category=Other bin=0 n
\newcommand{\CatOtherBinZeroTenN}{\ensuremath{973}}
% q1a_bins_by_category.csv: category=Other bin=0 mean_price
\newcommand{\CatOtherBinZeroTenPrice}{\ensuremath{1.9}\%}
% q1a_bins_by_category.csv: category=Other bin=0 observed_rate
\newcommand{\CatOtherBinZeroTenObs}{\ensuremath{1.8}\%}
% q1a_bins_by_category.csv: category=Other bin=0 diff_obs_minus_price
\newcommand{\CatOtherBinZeroTenDiff}{\ensuremath{-0.07}\%}
% q1a_bins_by_category.csv: category=Other bin=0 p_bh
\newcommand{\CatOtherBinZeroTenBhP}{\ensuremath{1.000}}
% q1a_bins_by_category.csv: category=Other bin=1 n
\newcommand{\CatOtherBinTenTwentyN}{\ensuremath{135}}
% q1a_bins_by_category.csv: category=Other bin=1 mean_price
\newcommand{\CatOtherBinTenTwentyPrice}{\ensuremath{14.9}\%}
% q1a_bins_by_category.csv: category=Other bin=1 observed_rate
\newcommand{\CatOtherBinTenTwentyObs}{\ensuremath{9.6}\%}
% q1a_bins_by_category.csv: category=Other bin=1 diff_obs_minus_price
\newcommand{\CatOtherBinTenTwentyDiff}{\ensuremath{-5.27}\%}
% q1a_bins_by_category.csv: category=Other bin=1 p_bh
\newcommand{\CatOtherBinTenTwentyBhP}{\ensuremath{0.345}}
% q1a_bins_by_category.csv: category=Politics bin=0 n
\newcommand{\CatPoliticsBinZeroTenN}{\ensuremath{3{,}426}}
% q1a_bins_by_category.csv: category=Politics bin=0 mean_price
\newcommand{\CatPoliticsBinZeroTenPrice}{\ensuremath{2.0}\%}
% q1a_bins_by_category.csv: category=Politics bin=0 observed_rate
\newcommand{\CatPoliticsBinZeroTenObs}{\ensuremath{1.6}\%}
% q1a_bins_by_category.csv: category=Politics bin=0 diff_obs_minus_price
\newcommand{\CatPoliticsBinZeroTenDiff}{\ensuremath{-0.36}\%}
% q1a_bins_by_category.csv: category=Politics bin=0 p_bh
\newcommand{\CatPoliticsBinZeroTenBhP}{\ensuremath{0.494}}
% q1a_bins_by_category.csv: category=Politics bin=1 n
\newcommand{\CatPoliticsBinTenTwentyN}{\ensuremath{494}}
% q1a_bins_by_category.csv: category=Politics bin=1 mean_price
\newcommand{\CatPoliticsBinTenTwentyPrice}{\ensuremath{14.7}\%}
% q1a_bins_by_category.csv: category=Politics bin=1 observed_rate
\newcommand{\CatPoliticsBinTenTwentyObs}{\ensuremath{9.1}\%}
% q1a_bins_by_category.csv: category=Politics bin=1 diff_obs_minus_price
\newcommand{\CatPoliticsBinTenTwentyDiff}{\ensuremath{-5.54}\%}
% q1a_bins_by_category.csv: category=Politics bin=1 p_bh
\newcommand{\CatPoliticsBinTenTwentyBhP}{\ensuremath{0.017}}
% q1a_bins_by_category.csv: category=Price Action bin=0 n
\newcommand{\CatPriceActionBinZeroTenN}{\ensuremath{405}}
% q1a_bins_by_category.csv: category=Price Action bin=0 mean_price
\newcommand{\CatPriceActionBinZeroTenPrice}{\ensuremath{1.0}\%}
% q1a_bins_by_category.csv: category=Price Action bin=0 observed_rate
\newcommand{\CatPriceActionBinZeroTenObs}{\ensuremath{1.0}\%}
% q1a_bins_by_category.csv: category=Price Action bin=0 diff_obs_minus_price
\newcommand{\CatPriceActionBinZeroTenDiff}{\ensuremath{-0.02}\%}
% q1a_bins_by_category.csv: category=Price Action bin=0 p_bh
\newcommand{\CatPriceActionBinZeroTenBhP}{\ensuremath{1.000}}
% q1a_bins_by_category.csv: category=Price Action bin=1 n
\newcommand{\CatPriceActionBinTenTwentyN}{\ensuremath{24}}
% q1a_bins_by_category.csv: category=Price Action bin=1 mean_price
\newcommand{\CatPriceActionBinTenTwentyPrice}{\ensuremath{15.1}\%}
% q1a_bins_by_category.csv: category=Price Action bin=1 observed_rate
\newcommand{\CatPriceActionBinTenTwentyObs}{\ensuremath{12.5}\%}
% q1a_bins_by_category.csv: category=Price Action bin=1 diff_obs_minus_price
\newcommand{\CatPriceActionBinTenTwentyDiff}{\ensuremath{-2.56}\%}
% q1a_bins_by_category.csv: category=Price Action bin=1 p_bh
\newcommand{\CatPriceActionBinTenTwentyBhP}{\ensuremath{1.000}}
% q1a_bins_by_category.csv: category=Sci-Tech bin=0 n
\newcommand{\CatSciTechBinZeroTenN}{\ensuremath{455}}
% q1a_bins_by_category.csv: category=Sci-Tech bin=0 mean_price
\newcommand{\CatSciTechBinZeroTenPrice}{\ensuremath{1.3}\%}
% q1a_bins_by_category.csv: category=Sci-Tech bin=0 observed_rate
\newcommand{\CatSciTechBinZeroTenObs}{\ensuremath{1.5}\%}
% q1a_bins_by_category.csv: category=Sci-Tech bin=0 diff_obs_minus_price
\newcommand{\CatSciTechBinZeroTenDiff}{\ensuremath{0.19}\%}
% q1a_bins_by_category.csv: category=Sci-Tech bin=0 p_bh
\newcommand{\CatSciTechBinZeroTenBhP}{\ensuremath{0.973}}
% q1a_bins_by_category.csv: category=Sci-Tech bin=1 n
\newcommand{\CatSciTechBinTenTwentyN}{\ensuremath{28}}
% q1a_bins_by_category.csv: category=Sci-Tech bin=1 mean_price
\newcommand{\CatSciTechBinTenTwentyPrice}{\ensuremath{14.2}\%}
% q1a_bins_by_category.csv: category=Sci-Tech bin=1 observed_rate
\newcommand{\CatSciTechBinTenTwentyObs}{\ensuremath{17.9}\%}
% q1a_bins_by_category.csv: category=Sci-Tech bin=1 diff_obs_minus_price
\newcommand{\CatSciTechBinTenTwentyDiff}{\ensuremath{3.63}\%}
% q1a_bins_by_category.csv: category=Sci-Tech bin=1 p_bh
\newcommand{\CatSciTechBinTenTwentyBhP}{\ensuremath{0.873}}
% q1a_bins_by_category.csv: category=Sports bin=0 n
\newcommand{\CatSportsBinZeroTenN}{\ensuremath{3{,}561}}
% q1a_bins_by_category.csv: category=Sports bin=0 mean_price
\newcommand{\CatSportsBinZeroTenPrice}{\ensuremath{3.3}\%}
% q1a_bins_by_category.csv: category=Sports bin=0 observed_rate
\newcommand{\CatSportsBinZeroTenObs}{\ensuremath{2.6}\%}
% q1a_bins_by_category.csv: category=Sports bin=0 diff_obs_minus_price
\newcommand{\CatSportsBinZeroTenDiff}{\ensuremath{-0.74}\%}
% q1a_bins_by_category.csv: category=Sports bin=0 p_bh
\newcommand{\CatSportsBinZeroTenBhP}{\ensuremath{0.088}}
% q1a_bins_by_category.csv: category=Sports bin=1 n
\newcommand{\CatSportsBinTenTwentyN}{\ensuremath{1{,}398}}
% q1a_bins_by_category.csv: category=Sports bin=1 mean_price
\newcommand{\CatSportsBinTenTwentyPrice}{\ensuremath{15.0}\%}
% q1a_bins_by_category.csv: category=Sports bin=1 observed_rate
\newcommand{\CatSportsBinTenTwentyObs}{\ensuremath{15.0}\%}
% q1a_bins_by_category.csv: category=Sports bin=1 diff_obs_minus_price
\newcommand{\CatSportsBinTenTwentyDiff}{\ensuremath{-0.03}\%}
% q1a_bins_by_category.csv: category=Sports bin=1 p_bh
\newcommand{\CatSportsBinTenTwentyBhP}{\ensuremath{1.000}}
% q1b_results.json: brier_diff_7d_minus_1d.n
\newcommand{\QOneBN}{\ensuremath{11{,}005}}
% q1b_results.json: brier_diff_7d_minus_1d.mean_diff
\newcommand{\QOneBMeanDiff}{\ensuremath{0.0257}}
% q1b_results.json: brier_diff_7d_minus_1d.ci_low
\newcommand{\QOneBCiLow}{\ensuremath{0.0232}}
% q1b_results.json: brier_diff_7d_minus_1d.ci_high
\newcommand{\QOneBCiHigh}{\ensuremath{0.0282}}
% q1b_results.json: brier_diff_7d_minus_1d.t
\newcommand{\QOneBT}{\ensuremath{20.0}}
% q1b_results.json: brier_diff_7d_minus_1d.df
\newcommand{\QOneBDf}{\ensuremath{11{,}004}}
% q1b_results.json: brier_diff_7d_minus_1d.p
\newcommand{\QOneBP}{\ensuremath{4.2\times10^{-87}}}
% q1b_results.json: brier_diff_7d_minus_1d.cohens_d
\newcommand{\QOneBD}{\ensuremath{0.19}}
% q1b_results.json: brier.7d
\newcommand{\QOneBBrierSevenDays}{\ensuremath{0.096}}
% q1b_results.json: brier.1d
\newcommand{\QOneBBrierOneDay}{\ensuremath{0.070}}
% q1b_results.json: abs_error_diff_7d_minus_1d.mean_diff
\newcommand{\QOneBAbsMeanDiff}{\ensuremath{0.0548}}
% q1b_results.json: abs_error_diff_7d_minus_1d.cohens_d
\newcommand{\QOneBAbsD}{\ensuremath{0.34}}
% q1b_results.json: mean_entropy_bits.diff_7d_minus_1d
\newcommand{\QOneBEntropyDiff}{\ensuremath{0.129}}
% q1b_results.json: ks_yes_vs_no.7d.D
\newcommand{\QOneBKsDSevenDays}{\ensuremath{0.705}}
% q1b_results.json: ks_yes_vs_no.1d.D
\newcommand{\QOneBKsDOneDay}{\ensuremath{0.786}}
% q1b_results.json: selection.max_staleness_h
\newcommand{\QOneBMaxStaleHours}{\ensuremath{24}}
% q1c_results.json: spearman_log_volume_vs_error.n
\newcommand{\QOneCVolumeN}{\ensuremath{47{,}948}}
% q1c_results.json: spearman_log_volume_vs_error.rho
\newcommand{\QOneCVolumeRho}{\ensuremath{0.226}}
% q1c_results.json: spearman_log_volume_vs_error.ci_low
\newcommand{\QOneCVolumeCiLow}{\ensuremath{0.217}}
% q1c_results.json: spearman_log_volume_vs_error.ci_high
\newcommand{\QOneCVolumeCiHigh}{\ensuremath{0.235}}
% q1c_results.json: spearman_log_volume_vs_error.p
\newcommand{\QOneCVolumeP}{\ensuremath{<10^{-300}}}
% q1c_results.json: spearman_log_trades_vs_error.n
\newcommand{\QOneCTradesN}{\ensuremath{47{,}948}}
% q1c_results.json: spearman_log_trades_vs_error.rho
\newcommand{\QOneCTradesRho}{\ensuremath{0.208}}
% q1c_results.json: spearman_log_trades_vs_error.ci_low
\newcommand{\QOneCTradesCiLow}{\ensuremath{0.199}}
% q1c_results.json: spearman_log_trades_vs_error.ci_high
\newcommand{\QOneCTradesCiHigh}{\ensuremath{0.217}}
% q1c_results.json: spearman_log_trades_vs_error.p
\newcommand{\QOneCTradesP}{\ensuremath{<10^{-300}}}
% q2_results.json: concentration.maker.n_wallets
\newcommand{\MakerNWallets}{\ensuremath{581{,}435}}
% q2_results.json: concentration.maker.gini
\newcommand{\MakerGini}{\ensuremath{0.978}}
% q2_results.json: concentration.maker.top1_share
\newcommand{\MakerTopOneShare}{\ensuremath{86.1}\%}
% q2_results.json: concentration.maker.top5_share
\newcommand{\MakerTopFiveShare}{\ensuremath{94.6}\%}
% q2_results.json: concentration.maker.top10_share
\newcommand{\MakerTopTenShare}{\ensuremath{96.9}\%}
% q2_results.json: concentration.maker.mean
\newcommand{\MakerMeanVolume}{\ensuremath{14{,}103}}
% q2_results.json: concentration.maker.median
\newcommand{\MakerMedianVolume}{\ensuremath{101.2}}
% q2_results.json: concentration.taker.n_wallets
\newcommand{\TakerNWallets}{\ensuremath{1{,}535{,}977}}
% q2_results.json: concentration.taker.gini
\newcommand{\TakerGini}{\ensuremath{0.963}}
% q2_results.json: concentration.taker.top1_share
\newcommand{\TakerTopOneShare}{\ensuremath{75.9}\%}
% q2_results.json: concentration.taker.top5_share
\newcommand{\TakerTopFiveShare}{\ensuremath{90.4}\%}
% q2_results.json: concentration.taker.top10_share
\newcommand{\TakerTopTenShare}{\ensuremath{94.8}\%}
% q2_results.json: concentration.taker.mean
\newcommand{\TakerMeanVolume}{\ensuremath{5{,}339}}
% q2_results.json: concentration.taker.median
\newcommand{\TakerMedianVolume}{\ensuremath{86.6}}
% q2_results.json: role_tests.welch_t.t
\newcommand{\RoleWelchT}{\ensuremath{117.1}}
% q2_results.json: role_tests.welch_t.p
\newcommand{\RoleWelchP}{\ensuremath{<10^{-300}}}
% q2_results.json: role_tests.welch_t.cohens_d
\newcommand{\RoleCohenD}{\ensuremath{0.18}}
% q2_results.json: role_tests.ks.D
\newcommand{\RoleKsD}{\ensuremath{0.057}}
% q2_results.json: spearman_n_trades_vs_volume.rho
\newcommand{\SpearmanTradesVolumeRho}{\ensuremath{0.655}}
% q2_results.json: kmeans.anova_mean_trade_price.eta_squared
\newcommand{\KmeansPriceEtaSq}{\ensuremath{0.008}}
% q2_results.json: kmeans.anova_mean_trade_price.perm_p
\newcommand{\KmeansPricePermP}{\ensuremath{1.0\times10^{-3}}}
% q2_results.json: kmeans.k_chosen
\newcommand{\KmeansK}{\ensuremath{4}}
% q2_results.json: kmeans.n_sample
\newcommand{\KmeansNSample}{\ensuremath{200{,}000}}
% q2_results.json: kmeans.robust_scaling_check.adjusted_rand_index_vs_zscore
\newcommand{\KmeansAri}{\ensuremath{0.93}}
% q2_results.json: kmeans.anova_mean_trade_price.n_perm
\newcommand{\KmeansNPerm}{\ensuremath{1{,}000}}
% q2_results.json: kmeans.chi2_cluster_x_category.cramers_v
\newcommand{\KmeansCategoryV}{\ensuremath{0.105}}
% q2_results.json: kmeans.cluster_summary[cluster=0].n
\newcommand{\ClusterZeroN}{\ensuremath{57{,}550}}
% q2_results.json: kmeans.cluster_summary[cluster=0].median_volume
\newcommand{\ClusterZeroMedianVolume}{\ensuremath{212.7}}
% q2_results.json: kmeans.cluster_summary[cluster=0].median_n_trades
\newcommand{\ClusterZeroMedianTrades}{\ensuremath{6}}
% q2_results.json: kmeans.cluster_summary[cluster=0].median_n_markets
\newcommand{\ClusterZeroMedianMarkets}{\ensuremath{2}}
% q2_results.json: kmeans.cluster_summary[cluster=1].n
\newcommand{\ClusterOneN}{\ensuremath{15{,}294}}
% q2_results.json: kmeans.cluster_summary[cluster=1].median_volume
\newcommand{\ClusterOneMedianVolume}{\ensuremath{5815.3}}
% q2_results.json: kmeans.cluster_summary[cluster=1].median_n_trades
\newcommand{\ClusterOneMedianTrades}{\ensuremath{129}}
% q2_results.json: kmeans.cluster_summary[cluster=1].median_n_markets
\newcommand{\ClusterOneMedianMarkets}{\ensuremath{42}}
% q2_results.json: kmeans.cluster_summary[cluster=2].n
\newcommand{\ClusterTwoN}{\ensuremath{69{,}556}}
% q2_results.json: kmeans.cluster_summary[cluster=2].median_volume
\newcommand{\ClusterTwoMedianVolume}{\ensuremath{6.6}}
% q2_results.json: kmeans.cluster_summary[cluster=2].median_n_trades
\newcommand{\ClusterTwoMedianTrades}{\ensuremath{2}}
% q2_results.json: kmeans.cluster_summary[cluster=2].median_n_markets
\newcommand{\ClusterTwoMedianMarkets}{\ensuremath{1}}
% q2_results.json: kmeans.cluster_summary[cluster=3].n
\newcommand{\ClusterThreeN}{\ensuremath{57{,}600}}
% q2_results.json: kmeans.cluster_summary[cluster=3].median_volume
\newcommand{\ClusterThreeMedianVolume}{\ensuremath{269.0}}
% q2_results.json: kmeans.cluster_summary[cluster=3].median_n_trades
\newcommand{\ClusterThreeMedianTrades}{\ensuremath{20}}
% q2_results.json: kmeans.cluster_summary[cluster=3].median_n_markets
\newcommand{\ClusterThreeMedianMarkets}{\ensuremath{10}}
% q3_results.json: summary.p_1d.mean
\newcommand{\PriceOneDayMean}{\ensuremath{0.376}}
% q3_results.json: summary.p_1d.median
\newcommand{\PriceOneDayMedian}{\ensuremath{0.320}}
% q3_results.json: summary.trade_size_usdc.mean
\newcommand{\TradeSizeMean}{\ensuremath{129.8}}
% q3_results.json: summary.trade_size_usdc.median
\newcommand{\TradeSizeMedian}{\ensuremath{6.93}}
% q3_results.json: fits.trade_size_usdc.power_law.alpha
\newcommand{\TradeSizePowerAlpha}{\ensuremath{1.88}}
% q3_results.json: fits.trade_size_usdc.power_law.xmin
\newcommand{\TradeSizePowerXmin}{\ensuremath{154.0}}
% q3_results.json: fits.trade_size_usdc.power_law.D
\newcommand{\TradeSizePowerKsD}{\ensuremath{0.039}}
% q3_results.json: fits.trade_size_usdc.exponential.D
\newcommand{\TradeSizeExponentialKsD}{\ensuremath{0.534}}
% q3_results.json: fits.trade_size_usdc.normal.D
\newcommand{\TradeSizeNormalKsD}{\ensuremath{0.469}}
% q3_results.json: fits.wallet_volume_usdc.power_law.alpha
\newcommand{\WalletVolumePowerAlpha}{\ensuremath{1.79}}
% q3_results.json: fits.wallet_volume_usdc.power_law.xmin
\newcommand{\WalletVolumePowerXmin}{\ensuremath{3450.2}}
% q3_results.json: fits.wallet_volume_usdc.power_law.D
\newcommand{\WalletVolumePowerKsD}{\ensuremath{0.024}}
% q3_results.json: fits.wallet_volume_usdc.exponential.D
\newcommand{\WalletVolumeExponentialKsD}{\ensuremath{0.675}}
% q3_results.json: fits.wallet_volume_usdc.normal.D
\newcommand{\WalletVolumeNormalKsD}{\ensuremath{0.488}}
% q1_bins_eqfreq.csv: row count
\newcommand{\QOneEqFreqBins}{\ensuremath{10}}
% q1_bins_eqfreq.csv: COMPUTED by this script (BH on the saved binom_p at verdicts.json alpha); the file has no p_bh column
\newcommand{\QOneEqFreqBinsBh}{\ensuremath{7}}
% close_time_summary.json (pilot run): n_markets
\newcommand{\PilotNMarkets}{\ensuremath{2{,}063}}
% close_time_summary.json (pilot run): share_markets_trade_after_close
\newcommand{\PilotShareMarketsAfterClose}{\ensuremath{51.3}\%}
% close_time_summary.json (pilot run): share_volume_after_close
\newcommand{\PilotShareVolumeAfterClose}{\ensuremath{34.0}\%}
% robustness run directory name
\newcommand{\RobustRun}{142123\_robust\_full\_v4}
% manifest_r1_split_payout.json: code_sha (hash of the script r1_split_payout.py alone)
\newcommand{\RobustCodeShaROne}{1232fa0fe1e2}
% manifest_r2_price_path.json: code_sha (hash of the script r2_price_path.py alone)
\newcommand{\RobustCodeShaRTwo}{4bcaceda3b5a}
% manifest_r1_split_payout.json, manifest_r2_price_path.json: code_version
\newcommand{\RobustCodeVersion}{v4}
% code.sha256: sha256 of the file that lists the hash of every script (first 12 hex)
\newcommand{\RobustCodeSha}{f91fc1970be1}
% path of the per-script hash list
\newcommand{\RobustCodeShaFile}{runs/142123\_robust\_full\_v4/code.sha256}
% r1_results.json: n_kept
\newcommand{\RobROneNKept}{\ensuremath{47{,}948}}
% r1_results.json: class_counts.clean
\newcommand{\RobROneCleanN}{\ensuremath{47{,}940}}
% r1_results.json: class_shares.clean
\newcommand{\RobROneCleanShare}{\ensuremath{99.98}\%}
% r1_results.json: class_counts.split
\newcommand{\RobROneSplitN}{\ensuremath{7}}
% r1_results.json: class_shares.split
\newcommand{\RobROneSplitShare}{\ensuremath{0.01}\%}
% r1_results.json: class_counts.other
\newcommand{\RobROneOtherN}{\ensuremath{0}}
% r1_results.json: class_shares.other
\newcommand{\RobROneOtherShare}{\ensuremath{0.00}\%}
% r1_results.json: class_counts.missing
\newcommand{\RobROneMissingN}{\ensuremath{1}}
% r1_results.json: class_shares.missing
\newcommand{\RobROneMissingShare}{\ensuremath{0.00}\%}
% r1_results.json: duplicates.kept_ids_with_duplicate_rows
\newcommand{\RobROneDupIds}{\ensuremath{0}}
% r1_results.json: duplicates.kept_duplicates_conflicting
\newcommand{\RobROneDupConflicting}{\ensuremath{0}}
% r1_results.json: our_y_on_split_markets.y1
\newcommand{\RobROneSplitYesOne}{\ensuremath{7}}
% r1_results.json: our_y_on_split_markets.y0
\newcommand{\RobROneSplitYesZero}{\ensuremath{0}}
% r1_results.json: mapping.n_clean
\newcommand{\RobROneNClean}{\ensuremath{47{,}940}}
% r1_results.json: mapping.agreement_index_primary
\newcommand{\RobROneAgreement}{\ensuremath{99.92}\%}
% r1_results.json: mapping.agreement_ci95[0]
\newcommand{\RobROneAgreementLow}{\ensuremath{99.89}\%}
% r1_results.json: mapping.agreement_ci95[1]
\newcommand{\RobROneAgreementHigh}{\ensuremath{99.94}\%}
% r1_results.json: mapping.n_mismatch
\newcommand{\RobROneNMismatch}{\ensuremath{40}}
% r1_results.json: mapping.validated
\newcommand{\RobROneMappingValidated}{yes}
% r1_results.json: sensitivity.n_excluded
\newcommand{\RobROneNExcluded}{\ensuremath{47}}
% r1_results.json: sensitivity.n_markets
\newcommand{\RobROneSensN}{\ensuremath{47{,}901}}
% r1_results.json: sensitivity.q1_sensitivity.chi2_gof.chi2
\newcommand{\RobROneSensChiSq}{\ensuremath{64.3}}
% r1_results.json: sensitivity.q1_sensitivity.chi2_gof.df
\newcommand{\RobROneSensDf}{\ensuremath{10}}
% r1_results.json: sensitivity.q1_sensitivity.chi2_gof.p
\newcommand{\RobROneSensP}{\ensuremath{5.6\times10^{-10}}}
% r1_results.json: sensitivity.q1_sensitivity.chi2_gof.effect_w
\newcommand{\RobROneSensW}{\ensuremath{0.037}}
% r1_results.json: sensitivity.q1_sensitivity.brier
\newcommand{\RobROneSensBrier}{\ensuremath{0.107}}
% r1_results.json: sensitivity.q1_sensitivity.bins_reject_bh
\newcommand{\RobROneSensBinsBh}{\ensuremath{4}}
% r1_results.json: sensitivity.delta_effect_w
\newcommand{\RobROneDeltaW}{\ensuremath{-0.0010}}
% r1_results.json: sensitivity.rule_results_same
\newcommand{\RobROneRuleSame}{yes}
% r1_results.json: sensitivity.rule_result_original
\newcommand{\RobROneRuleOriginal}{Prices are not calibrated as a whole at this n.}
% r1_results.json: sensitivity.rule_result_sensitivity
\newcommand{\RobROneRuleSens}{Prices are not calibrated as a whole at this n.}
% r2_results.json: overall.n
\newcommand{\RobRTwoN}{\ensuremath{47{,}948}}
% r2_results.json: overall.median_abs_delta
\newcommand{\RobRTwoMedianAbsDelta}{\ensuremath{0.0153}}
% r2_results.json: overall.share_no_trade
\newcommand{\RobRTwoShareNoTrade}{\ensuremath{5.9}\%}
% r2_results.json: overall.share_crossing
\newcommand{\RobRTwoShareCrossing}{\ensuremath{8.02}\%}
% r2_results.json: overall.share_abs_delta_ge_0.10
\newcommand{\RobRTwoShareDeltaAtLeastTen}{\ensuremath{18.98}\%}
% r2_results.json: overall.share_abs_delta_ge_0.25
\newcommand{\RobRTwoShareDeltaAtLeastTwentyFive}{\ensuremath{10.97}\%}
% r2_results.json: overall.share_abs_delta_ge_0.50
\newcommand{\RobRTwoShareDeltaAtLeastFifty}{\ensuremath{4.65}\%}
% r2_results.json: overall.quantiles_abs_delta.q0.5
\newcommand{\RobRTwoAbsDeltaMedian}{\ensuremath{0.015}}
% r2_results.json: overall.quantiles_abs_delta.q0.75
\newcommand{\RobRTwoAbsDeltaSevenFive}{\ensuremath{0.060}}
% r2_results.json: overall.quantiles_abs_delta.q0.9
\newcommand{\RobRTwoAbsDeltaNinety}{\ensuremath{0.280}}
% r2_results.json: overall.quantiles_abs_delta.q0.95
\newcommand{\RobRTwoAbsDeltaNinetyFive}{\ensuremath{0.482}}
% r2_results.json: overall.quantiles_abs_delta.q0.99
\newcommand{\RobRTwoAbsDeltaNinetyNine}{\ensuremath{0.789}}
% r2_results.json: overall.quantiles_abs_delta.max
\newcommand{\RobRTwoAbsDeltaMax}{\ensuremath{0.998}}
% r2_results.json: overall.n_window_trades.median
\newcommand{\RobRTwoWindowTradesMedian}{\ensuremath{61}}
% r2_results.json: overall.n_window_trades.q0.9
\newcommand{\RobRTwoWindowTradesNinety}{\ensuremath{1{,}147}}
% r2_results.json: overall.n_window_trades.max
\newcommand{\RobRTwoWindowTradesMax}{\ensuremath{40{,}499}}
% r2_results.json: overall.brier.n
\newcommand{\RobRTwoAllN}{\ensuremath{47{,}948}}
% r2_results.json: overall.brier.brier_1d
\newcommand{\RobRTwoAllBrierOneDay}{\ensuremath{0.1071}}
% r2_results.json: overall.brier.brier_T
\newcommand{\RobRTwoAllBrierAtClose}{\ensuremath{0.0749}}
% r2_results.json: overall.brier.mean_diff
\newcommand{\RobRTwoAllMeanDiff}{\ensuremath{-0.0322}}
% r2_results.json: overall.brier.ci_low
\newcommand{\RobRTwoAllCiLow}{\ensuremath{-0.0333}}
% r2_results.json: overall.brier.ci_high
\newcommand{\RobRTwoAllCiHigh}{\ensuremath{-0.0311}}
% r2_results.json: overall.brier.t
\newcommand{\RobRTwoAllT}{\ensuremath{-56.5}}
% r2_results.json: overall.brier.p
\newcommand{\RobRTwoAllP}{\ensuremath{<10^{-300}}}
% r2_results.json: overall.brier.cohens_dz
\newcommand{\RobRTwoAllDz}{\ensuremath{-0.26}}
% r2_results.json: overall.brier_traded_only.n
\newcommand{\RobRTwoTradedN}{\ensuremath{45{,}112}}
% r2_results.json: overall.brier_traded_only.brier_1d
\newcommand{\RobRTwoTradedBrierOneDay}{\ensuremath{0.1122}}
% r2_results.json: overall.brier_traded_only.brier_T
\newcommand{\RobRTwoTradedBrierAtClose}{\ensuremath{0.0780}}
% r2_results.json: overall.brier_traded_only.mean_diff
\newcommand{\RobRTwoTradedMeanDiff}{\ensuremath{-0.0342}}
% r2_results.json: overall.brier_traded_only.ci_low
\newcommand{\RobRTwoTradedCiLow}{\ensuremath{-0.0354}}
% r2_results.json: overall.brier_traded_only.ci_high
\newcommand{\RobRTwoTradedCiHigh}{\ensuremath{-0.0330}}
% r2_results.json: overall.brier_traded_only.t
\newcommand{\RobRTwoTradedT}{\ensuremath{-56.6}}
% r2_results.json: overall.brier_traded_only.p
\newcommand{\RobRTwoTradedP}{\ensuremath{<10^{-300}}}
% r2_results.json: overall.brier_traded_only.cohens_dz
\newcommand{\RobRTwoTradedDz}{\ensuremath{-0.27}}
% r2_results.json: by_category.Crypto.n
\newcommand{\RobRTwoCatCryptoN}{\ensuremath{11{,}888}}
% r2_results.json: by_category.Crypto.median_abs_delta
\newcommand{\RobRTwoCatCryptoMedianAbsDelta}{\ensuremath{0.0040}}
% r2_results.json: by_category.Crypto.brier.mean_diff
\newcommand{\RobRTwoCatCryptoMeanDiff}{\ensuremath{-0.0322}}
% r2_results.json: by_category.Culture.n
\newcommand{\RobRTwoCatCultureN}{\ensuremath{1{,}074}}
% r2_results.json: by_category.Culture.median_abs_delta
\newcommand{\RobRTwoCatCultureMedianAbsDelta}{\ensuremath{0.0100}}
% r2_results.json: by_category.Culture.brier.mean_diff
\newcommand{\RobRTwoCatCultureMeanDiff}{\ensuremath{-0.0083}}
% r2_results.json: by_category.Finance.n
\newcommand{\RobRTwoCatFinanceN}{\ensuremath{5{,}146}}
% r2_results.json: by_category.Finance.median_abs_delta
\newcommand{\RobRTwoCatFinanceMedianAbsDelta}{\ensuremath{0.0150}}
% r2_results.json: by_category.Finance.brier.mean_diff
\newcommand{\RobRTwoCatFinanceMeanDiff}{\ensuremath{-0.0424}}
% r2_results.json: by_category.Other.n
\newcommand{\RobRTwoCatOtherN}{\ensuremath{2{,}274}}
% r2_results.json: by_category.Other.median_abs_delta
\newcommand{\RobRTwoCatOtherMedianAbsDelta}{\ensuremath{0.0110}}
% r2_results.json: by_category.Other.brier.mean_diff
\newcommand{\RobRTwoCatOtherMeanDiff}{\ensuremath{-0.0238}}
% r2_results.json: by_category.Politics.n
\newcommand{\RobRTwoCatPoliticsN}{\ensuremath{6{,}922}}
% r2_results.json: by_category.Politics.median_abs_delta
\newcommand{\RobRTwoCatPoliticsMedianAbsDelta}{\ensuremath{0.0125}}
% r2_results.json: by_category.Politics.brier.mean_diff
\newcommand{\RobRTwoCatPoliticsMeanDiff}{\ensuremath{-0.0140}}
% r2_results.json: by_category.Price Action.n
\newcommand{\RobRTwoCatPriceActionN}{\ensuremath{663}}
% r2_results.json: by_category.Price Action.median_abs_delta
\newcommand{\RobRTwoCatPriceActionMedianAbsDelta}{\ensuremath{0.0030}}
% r2_results.json: by_category.Price Action.brier.mean_diff
\newcommand{\RobRTwoCatPriceActionMeanDiff}{\ensuremath{-0.0470}}
% r2_results.json: by_category.Sci-Tech.n
\newcommand{\RobRTwoCatSciTechN}{\ensuremath{819}}
% r2_results.json: by_category.Sci-Tech.median_abs_delta
\newcommand{\RobRTwoCatSciTechMedianAbsDelta}{\ensuremath{0.0073}}
% r2_results.json: by_category.Sci-Tech.brier.mean_diff
\newcommand{\RobRTwoCatSciTechMeanDiff}{\ensuremath{-0.0130}}
% r2_results.json: by_category.Sports.n
\newcommand{\RobRTwoCatSportsN}{\ensuremath{19{,}162}}
% r2_results.json: by_category.Sports.median_abs_delta
\newcommand{\RobRTwoCatSportsMedianAbsDelta}{\ensuremath{0.0200}}
% r2_results.json: by_category.Sports.brier.mean_diff
\newcommand{\RobRTwoCatSportsMeanDiff}{\ensuremath{-0.0386}}
% 142123_robust_full_v4/markets.csv: category_refined=Crypto, row count
\newcommand{\CatCryptoNInScope}{\ensuremath{102{,}204}}
% 142123_robust_full_v4/markets.csv: category_refined=Crypto, row count / all rows
\newcommand{\CatCryptoShareInScope}{\ensuremath{14.9}\%}
% 142123_robust_full_v4/markets.csv: category_refined=Crypto, kept=True
\newcommand{\CatCryptoNKept}{\ensuremath{11{,}888}}
% 142123_robust_full_v4/markets.csv: category_refined=Crypto, kept=True / all kept
\newcommand{\CatCryptoShareKept}{\ensuremath{24.8}\%}
% 142123_robust_full_v4/markets.csv: category_refined=Crypto, kept=True / row count
\newcommand{\CatCryptoSurvival}{\ensuremath{11.6}\%}
% 142123_robust_full_v4/markets.csv: category_refined=Crypto, drop_reason=updown
\newcommand{\CatCryptoDropUpdownN}{\ensuremath{60{,}742}}
% 142123_robust_full_v4/markets.csv: category_refined=Crypto, drop_reason=no_outcome
\newcommand{\CatCryptoDropNoOutcomeN}{\ensuremath{7{,}696}}
% 142123_robust_full_v4/markets.csv: category_refined=Crypto, drop_reason=no_price
\newcommand{\CatCryptoDropNoPriceN}{\ensuremath{18{,}433}}
% 142123_robust_full_v4/markets.csv: category_refined=Crypto, drop_reason=stale_price
\newcommand{\CatCryptoDropStalePriceN}{\ensuremath{2{,}721}}
% 142123_robust_full_v4/markets.csv: category_refined=Crypto, drop_reason=below_min_trades
\newcommand{\CatCryptoDropBelowMinTradesN}{\ensuremath{724}}
% 142123_robust_full_v4/markets.csv: category_refined=Culture, row count
\newcommand{\CatCultureNInScope}{\ensuremath{2{,}999}}
% 142123_robust_full_v4/markets.csv: category_refined=Culture, row count / all rows
\newcommand{\CatCultureShareInScope}{\ensuremath{0.4}\%}
% 142123_robust_full_v4/markets.csv: category_refined=Culture, kept=True
\newcommand{\CatCultureNKept}{\ensuremath{1{,}074}}
% 142123_robust_full_v4/markets.csv: category_refined=Culture, kept=True / all kept
\newcommand{\CatCultureShareKept}{\ensuremath{2.2}\%}
% 142123_robust_full_v4/markets.csv: category_refined=Culture, kept=True / row count
\newcommand{\CatCultureSurvival}{\ensuremath{35.8}\%}
% 142123_robust_full_v4/markets.csv: category_refined=Culture, drop_reason=updown
\newcommand{\CatCultureDropUpdownN}{\ensuremath{0}}
% 142123_robust_full_v4/markets.csv: category_refined=Culture, drop_reason=no_outcome
\newcommand{\CatCultureDropNoOutcomeN}{\ensuremath{280}}
% 142123_robust_full_v4/markets.csv: category_refined=Culture, drop_reason=no_price
\newcommand{\CatCultureDropNoPriceN}{\ensuremath{368}}
% 142123_robust_full_v4/markets.csv: category_refined=Culture, drop_reason=stale_price
\newcommand{\CatCultureDropStalePriceN}{\ensuremath{1{,}140}}
% 142123_robust_full_v4/markets.csv: category_refined=Culture, drop_reason=below_min_trades
\newcommand{\CatCultureDropBelowMinTradesN}{\ensuremath{137}}
% 142123_robust_full_v4/markets.csv: category_refined=Finance, row count
\newcommand{\CatFinanceNInScope}{\ensuremath{13{,}370}}
% 142123_robust_full_v4/markets.csv: category_refined=Finance, row count / all rows
\newcommand{\CatFinanceShareInScope}{\ensuremath{1.9}\%}
% 142123_robust_full_v4/markets.csv: category_refined=Finance, kept=True
\newcommand{\CatFinanceNKept}{\ensuremath{5{,}146}}
% 142123_robust_full_v4/markets.csv: category_refined=Finance, kept=True / all kept
\newcommand{\CatFinanceShareKept}{\ensuremath{10.7}\%}
% 142123_robust_full_v4/markets.csv: category_refined=Finance, kept=True / row count
\newcommand{\CatFinanceSurvival}{\ensuremath{38.5}\%}
% 142123_robust_full_v4/markets.csv: category_refined=Finance, drop_reason=updown
\newcommand{\CatFinanceDropUpdownN}{\ensuremath{1{,}574}}
% 142123_robust_full_v4/markets.csv: category_refined=Finance, drop_reason=no_outcome
\newcommand{\CatFinanceDropNoOutcomeN}{\ensuremath{996}}
% 142123_robust_full_v4/markets.csv: category_refined=Finance, drop_reason=no_price
\newcommand{\CatFinanceDropNoPriceN}{\ensuremath{1{,}117}}
% 142123_robust_full_v4/markets.csv: category_refined=Finance, drop_reason=stale_price
\newcommand{\CatFinanceDropStalePriceN}{\ensuremath{2{,}506}}
% 142123_robust_full_v4/markets.csv: category_refined=Finance, drop_reason=below_min_trades
\newcommand{\CatFinanceDropBelowMinTradesN}{\ensuremath{2{,}031}}
% 142123_robust_full_v4/markets.csv: category_refined=Other, row count
\newcommand{\CatOtherNInScope}{\ensuremath{13{,}775}}
% 142123_robust_full_v4/markets.csv: category_refined=Other, row count / all rows
\newcommand{\CatOtherShareInScope}{\ensuremath{2.0}\%}
% 142123_robust_full_v4/markets.csv: category_refined=Other, kept=True
\newcommand{\CatOtherNKept}{\ensuremath{2{,}274}}
% 142123_robust_full_v4/markets.csv: category_refined=Other, kept=True / all kept
\newcommand{\CatOtherShareKept}{\ensuremath{4.7}\%}
% 142123_robust_full_v4/markets.csv: category_refined=Other, kept=True / row count
\newcommand{\CatOtherSurvival}{\ensuremath{16.5}\%}
% 142123_robust_full_v4/markets.csv: category_refined=Other, drop_reason=updown
\newcommand{\CatOtherDropUpdownN}{\ensuremath{1{,}406}}
% 142123_robust_full_v4/markets.csv: category_refined=Other, drop_reason=no_outcome
\newcommand{\CatOtherDropNoOutcomeN}{\ensuremath{1{,}426}}
% 142123_robust_full_v4/markets.csv: category_refined=Other, drop_reason=no_price
\newcommand{\CatOtherDropNoPriceN}{\ensuremath{5{,}493}}
% 142123_robust_full_v4/markets.csv: category_refined=Other, drop_reason=stale_price
\newcommand{\CatOtherDropStalePriceN}{\ensuremath{2{,}044}}
% 142123_robust_full_v4/markets.csv: category_refined=Other, drop_reason=below_min_trades
\newcommand{\CatOtherDropBelowMinTradesN}{\ensuremath{1{,}132}}
% 142123_robust_full_v4/markets.csv: category_refined=Politics, row count
\newcommand{\CatPoliticsNInScope}{\ensuremath{18{,}250}}
% 142123_robust_full_v4/markets.csv: category_refined=Politics, row count / all rows
\newcommand{\CatPoliticsShareInScope}{\ensuremath{2.7}\%}
% 142123_robust_full_v4/markets.csv: category_refined=Politics, kept=True
\newcommand{\CatPoliticsNKept}{\ensuremath{6{,}922}}
% 142123_robust_full_v4/markets.csv: category_refined=Politics, kept=True / all kept
\newcommand{\CatPoliticsShareKept}{\ensuremath{14.4}\%}
% 142123_robust_full_v4/markets.csv: category_refined=Politics, kept=True / row count
\newcommand{\CatPoliticsSurvival}{\ensuremath{37.9}\%}
% 142123_robust_full_v4/markets.csv: category_refined=Politics, drop_reason=updown
\newcommand{\CatPoliticsDropUpdownN}{\ensuremath{56}}
% 142123_robust_full_v4/markets.csv: category_refined=Politics, drop_reason=no_outcome
\newcommand{\CatPoliticsDropNoOutcomeN}{\ensuremath{1{,}527}}
% 142123_robust_full_v4/markets.csv: category_refined=Politics, drop_reason=no_price
\newcommand{\CatPoliticsDropNoPriceN}{\ensuremath{4{,}973}}
% 142123_robust_full_v4/markets.csv: category_refined=Politics, drop_reason=stale_price
\newcommand{\CatPoliticsDropStalePriceN}{\ensuremath{4{,}035}}
% 142123_robust_full_v4/markets.csv: category_refined=Politics, drop_reason=below_min_trades
\newcommand{\CatPoliticsDropBelowMinTradesN}{\ensuremath{737}}
% 142123_robust_full_v4/markets.csv: category_refined=Price Action, row count
\newcommand{\CatPriceActionNInScope}{\ensuremath{192{,}762}}
% 142123_robust_full_v4/markets.csv: category_refined=Price Action, row count / all rows
\newcommand{\CatPriceActionShareInScope}{\ensuremath{28.0}\%}
% 142123_robust_full_v4/markets.csv: category_refined=Price Action, kept=True
\newcommand{\CatPriceActionNKept}{\ensuremath{663}}
% 142123_robust_full_v4/markets.csv: category_refined=Price Action, kept=True / all kept
\newcommand{\CatPriceActionShareKept}{\ensuremath{1.4}\%}
% 142123_robust_full_v4/markets.csv: category_refined=Price Action, kept=True / row count
\newcommand{\CatPriceActionSurvival}{\ensuremath{0.3}\%}
% 142123_robust_full_v4/markets.csv: category_refined=Price Action, drop_reason=updown
\newcommand{\CatPriceActionDropUpdownN}{\ensuremath{191{,}290}}
% 142123_robust_full_v4/markets.csv: category_refined=Price Action, drop_reason=no_outcome
\newcommand{\CatPriceActionDropNoOutcomeN}{\ensuremath{214}}
% 142123_robust_full_v4/markets.csv: category_refined=Price Action, drop_reason=no_price
\newcommand{\CatPriceActionDropNoPriceN}{\ensuremath{373}}
% 142123_robust_full_v4/markets.csv: category_refined=Price Action, drop_reason=stale_price
\newcommand{\CatPriceActionDropStalePriceN}{\ensuremath{183}}
% 142123_robust_full_v4/markets.csv: category_refined=Price Action, drop_reason=below_min_trades
\newcommand{\CatPriceActionDropBelowMinTradesN}{\ensuremath{39}}
% 142123_robust_full_v4/markets.csv: category_refined=Sci-Tech, row count
\newcommand{\CatSciTechNInScope}{\ensuremath{2{,}860}}
% 142123_robust_full_v4/markets.csv: category_refined=Sci-Tech, row count / all rows
\newcommand{\CatSciTechShareInScope}{\ensuremath{0.4}\%}
% 142123_robust_full_v4/markets.csv: category_refined=Sci-Tech, kept=True
\newcommand{\CatSciTechNKept}{\ensuremath{819}}
% 142123_robust_full_v4/markets.csv: category_refined=Sci-Tech, kept=True / all kept
\newcommand{\CatSciTechShareKept}{\ensuremath{1.7}\%}
% 142123_robust_full_v4/markets.csv: category_refined=Sci-Tech, kept=True / row count
\newcommand{\CatSciTechSurvival}{\ensuremath{28.6}\%}
% 142123_robust_full_v4/markets.csv: category_refined=Sci-Tech, drop_reason=updown
\newcommand{\CatSciTechDropUpdownN}{\ensuremath{233}}
% 142123_robust_full_v4/markets.csv: category_refined=Sci-Tech, drop_reason=no_outcome
\newcommand{\CatSciTechDropNoOutcomeN}{\ensuremath{252}}
% 142123_robust_full_v4/markets.csv: category_refined=Sci-Tech, drop_reason=no_price
\newcommand{\CatSciTechDropNoPriceN}{\ensuremath{533}}
% 142123_robust_full_v4/markets.csv: category_refined=Sci-Tech, drop_reason=stale_price
\newcommand{\CatSciTechDropStalePriceN}{\ensuremath{794}}
% 142123_robust_full_v4/markets.csv: category_refined=Sci-Tech, drop_reason=below_min_trades
\newcommand{\CatSciTechDropBelowMinTradesN}{\ensuremath{229}}
% 142123_robust_full_v4/markets.csv: category_refined=Sports, row count
\newcommand{\CatSportsNInScope}{\ensuremath{341{,}911}}
% 142123_robust_full_v4/markets.csv: category_refined=Sports, row count / all rows
\newcommand{\CatSportsShareInScope}{\ensuremath{49.7}\%}
% 142123_robust_full_v4/markets.csv: category_refined=Sports, kept=True
\newcommand{\CatSportsNKept}{\ensuremath{19{,}162}}
% 142123_robust_full_v4/markets.csv: category_refined=Sports, kept=True / all kept
\newcommand{\CatSportsShareKept}{\ensuremath{40.0}\%}
% 142123_robust_full_v4/markets.csv: category_refined=Sports, kept=True / row count
\newcommand{\CatSportsSurvival}{\ensuremath{5.6}\%}
% 142123_robust_full_v4/markets.csv: category_refined=Sports, row count
\newcommand{\SportsNInScope}{\ensuremath{341{,}911}}
% 142123_robust_full_v4/markets.csv: category_refined=Sports, row count / all rows
\newcommand{\SportsShareInScope}{\ensuremath{49.7}\%}
% 142123_robust_full_v4/markets.csv: category_refined=Sports, kept=True / all kept
\newcommand{\SportsShareKept}{\ensuremath{40.0}\%}
% 142123_robust_full_v4/markets.csv: category_refined=Sports, drop_reason=updown
\newcommand{\CatSportsDropUpdownN}{\ensuremath{119}}
% 142123_robust_full_v4/markets.csv: category_refined=Sports, drop_reason=no_outcome
\newcommand{\CatSportsDropNoOutcomeN}{\ensuremath{50{,}327}}
% 142123_robust_full_v4/markets.csv: category_refined=Sports, drop_reason=no_price
\newcommand{\CatSportsDropNoPriceN}{\ensuremath{183{,}999}}
% 142123_robust_full_v4/markets.csv: category_refined=Sports, drop_reason=stale_price
\newcommand{\CatSportsDropStalePriceN}{\ensuremath{66{,}807}}
% 142123_robust_full_v4/markets.csv: category_refined=Sports, drop_reason=below_min_trades
\newcommand{\CatSportsDropBelowMinTradesN}{\ensuremath{21{,}497}}


% --- Design constants and literature facts. Not measured by our runs. ---
% src: S-PAPER-3, S-PAPER-4 (design constants)
\newcommand{\NBins}{10}
\newcommand{\CiLevel}{95\%}
\newcommand{\DaysShort}{1}
\newcommand{\DaysLong}{7}
\newcommand{\MaxStaleHours}{24}
\newcommand{\TopOneLabel}{1\%}
\newcommand{\TopFiveLabel}{5\%}
\newcommand{\TopTenLabel}{10\%}
\newcommand{\HookPct}{70\%}
\newcommand{\pT}{\ensuremath{p_{T}}}
\newcommand{\DeltaA}{0.10}
\newcommand{\DeltaB}{0.25}
\newcommand{\DeltaC}{0.50}
\newcommand{\PayoutDollar}{\$1}
\newcommand{\HookPrice}{\$0.70}
\newcommand{\FavoriteLine}{0.5}
\newcommand{\pOne}{\ensuremath{p_{\DaysShort\mathrm{d}}}}
\newcommand{\pSeven}{\ensuremath{p_{\DaysLong\mathrm{d}}}}
\newcommand{\MinCategoryN}{30}
\newcommand{\MinBinN}{5}
\newcommand{\PermSeed}{42}
% src: S-PAPER-1, S-PAPER-2 (facts about the dataset paper)
\newcommand{\PaperTrades}{1.2 billion}
\newcommand{\PaperMarkets}{1.3 million}
\newcommand{\PaperVolume}{\$61 billion}
\newcommand{\PaperPeriod}{November 2022 to April 2026}
\newcommand{\PaperBrierMarkets}{1{,}019}
\newcommand{\PaperTopShare}{84.1\%}
\newcommand{\PaperGiniMaker}{0.970}
\newcommand{\PaperGiniTaker}{0.943}
\newcommand{\PaperCheapGap}{0.2 to 0.3 cents}
\newcommand{\PaperCheapPrice}{30 cents}

% --- Table row helpers: pull one row of macros by name ---
\newcommand{\binrow}[2]{#1 & \csname Bin#2N\endcsname & \csname Bin#2Price\endcsname
  & \csname Bin#2Obs\endcsname & [\csname Bin#2CiLow\endcsname, \csname Bin#2CiHigh\endcsname]
  & \csname Bin#2Diff\endcsname & \csname Bin#2BhP\endcsname}
\newcommand{\topicrow}[4]{#1 & \csname Cat#2N\endcsname & \csname Cat#2FavWin\endcsname
  & \csname Cat#2Brier\endcsname & \csname #3N\endcsname & \csname #3Diff\endcsname & \csname #3BhP\endcsname
  & \csname #4N\endcsname & \csname #4Diff\endcsname & \csname #4BhP\endcsname}

\title{When Polymarket Says \HookPct, Does It Happen \HookPct{} of the Time?\\[2pt]
\large Prediction Accuracy, Trading Concentration, and Statistical Distributions in Prediction Markets\\[4pt]
\normalsize Version 2: MVP1 + MVP2}
\author{Group 1 $|$ Team Brier: Taek Lim, Zhugang Liu (Tony), Francisco Cortez\\
CS 6330 Introduction to Data Science}
\date{}

\begin{document}
\maketitle

\begin{abstract}
% src: S-PAPER-1, S-PAPER-3, S-PAPER-5
A Polymarket YES share pays one dollar if the event happens.
So a price of \HookPct{} is read as a \HookPct{} chance.
We test this on resolved Yes/No markets from the Polymarket-v1 database.
We keep \NMarketsKept{} of \NMarketsInScope{} markets in scope.
We read each market's price \DaysShort{} day before its scheduled close.
The pre-registered chi-square test says the prices are not calibrated as a whole
($\chi^2=\QOneChiSq$, $p=\QOneP$, Cohen's $w=\QOneW$).
The observed Yes rate lies below the mean price in every price bin except the top one.
The top \TopOneLabel{} of maker wallets hold \MakerTopOneShare{} of maker volume
(Gini coefficient \MakerGini).
Accuracy depends on the topic (Cramér's $V=\QOneAV$).
Prices are closer to the outcome \DaysShort{} day before the close than \DaysLong{} days before.
The mean Brier difference is \QOneBMeanDiff{} (Cohen's $d=\QOneBD$).
We find no evidence that busier markets are more accurate (Spearman $\rho=\QOneCVolumeRho$).
Mean volume is far above median volume.
Trade sizes are strongly right skewed.
We describe Q2 and Q3 without a verdict.
A payout check leaves the Q1 verdict unchanged.
\end{abstract}

\section{Introduction}\label{sec:intro}
% src: S-PAPER-1
A Polymarket YES share pays \PayoutDollar{} if the event happens.
At a price of \HookPrice, the price is read as a \HookPct{} chance.
Is that chance right?
If we collect every event priced at \HookPct, about \HookPct{} of them should have happened.
The Polymarket-v1 database records every trade on Polymarket's first exchange.
It holds \PaperTrades{} trades and \PaperMarkets{} markets from \PaperPeriod{} \citep{qin2026polymarketv1}.
So we can check each price against what happened.

This paper reports two stages (MVP1 and MVP2) of a staged study.
% src: S-PAPER-8
The first stage answers three questions.
The first question is Q1: is the crowd right?
% src: S-PAPER-1
We test calibration one price range at a time.
We use binomial and chi-square tests, confidence intervals, and the Brier score.
We also ask whether cheap longshot bets are a trap.
The second question is Q2: who is the crowd?
We ask whether a small group of wallets does most of the trading.
We use Lorenz curves and the Gini coefficient.
We also group wallets by k-means, as a side analysis.
The third question is Q3: what does betting look like?
We describe trade sizes, wallet volumes, and prices with histograms and distribution fits.
% src: S-PAPER-2, S-PAPER-8
The second stage adds three sub-questions under Q1.
Q1a asks whether accuracy depends on the topic.
Q1b asks whether the market gets more accurate near the deadline.
Q1c asks whether busier markets are more accurate.
Each stage fixes its criteria before its results.

% src: S-PAPER-3, S-PAPER-5
Our contributions are these.
We fix the evaluation rules before we see any Q1 result.
We report one rule change made afterwards.
We score all \NMarketsKept{} kept markets, not a small panel.
We split the Q1 result by topic, by horizon, and by volume.
We report an effect size next to every p-value.
The text uses the word wallet for an address on the exchange.

\section{Related Work}\label{sec:related}
% src: S-PAPER-1
We build on two references.
Snowberg and Wolfers study horse racing \citep{snowberg2010favorite}.
They show that longshots win less often than their odds say.
We test the same gap on Polymarket.

The dataset paper reports Polymarket results on cheap tokens \citep[Table~3]{qin2026polymarketv1}.
Tokens priced under \PaperCheapPrice{} pay out \PaperCheapGap{} less than their price.
This is counted by trade.
We count markets, not trades.
We read this result as the same direction as the horse-racing bias.

The dataset paper also scores \PaperBrierMarkets{} markets with the Brier score \citep{qin2026polymarketv1}.
We score all resolved Yes/No markets that pass our filters.
Finally, the paper reports that the top \TopOneLabel{} of maker addresses hold \PaperTopShare{}
of maker volume \citep[Table~5]{qin2026polymarketv1}.
We ask whether a few wallets dominate volume in our kept markets.
% src: S-PAPER-2
\todo{owner: related work beyond the two references in the proposal, if the course report needs it}

\section{Data}\label{sec:data}
\subsection{Source and scope}\label{sec:data-source}
% src: S-PAPER-1, S-PAPER-2
We use the \texttt{daily\_aligned/} layer of Polymarket-v1.
This layer holds Standard Binary markets (\texttt{neg\_risk} is false).
Its relayer records are already removed \citep{polymarketv1data}.
We keep markets that have a winning outcome label.
We drop Up/Down markets by the default rule in our code.
We take the price from the \texttt{p\_event} column, which is on the YES scale.
The data are at \url{https://huggingface.co/datasets/TimeSeventeen/Polymarket-v1} \citep{polymarketv1data}.
The dataset card states the license CC-BY-4.0.
The paper that describes the data is \citep{qin2026polymarketv1}.

\subsection{Filters and drop reasons}\label{sec:data-filters}
% src: S-PAPER-5
In the full run, \NMarketsInScope{} markets are in scope.
We keep \NMarketsKept{} of them (\PctMarketsKept).
\Cref{fig:drops} counts each dropped market once, by the first rule it fails.
The rules apply in the order of the bars.
Up/Down markets are \NUpdown{} (\PctUpdown).
\NNoOutcome{} markets have no usable winner label (\PctNoOutcome).
\NNoPrice{} have no price at or before $T-\DaysShort$ day (\PctNoPrice).
\NStalePrice{} have a price older than \MaxStaleHours{} hours (\PctStalePrice).
\NBelowMinTrades{} have fewer than \MinTrades{} trades before $T-\DaysShort$ day (\PctBelowMinTrades).

\begin{figure}[H]
\centering
\figorbox{figures/drop_reasons}{0.65\linewidth}
\caption{Markets by first failing rule. Why: it shows which rule removes the most markets. Source: run \RunId.}
\label{fig:drops}
\end{figure}

\subsection{Market and wallet tables}\label{sec:data-tables}
% src: S-PAPER-2, S-PAPER-5
We build one market table and one wallet table.
The market table has one row per kept market.
It holds the category, trade count, volume in USDC, the prices \pOne{} and \pSeven, and the outcome $y\in\{0,1\}$.
The wallet table has one row per wallet and role (maker or taker).
It holds volume in USDC, trade count, market count, and mean trade price.
A wallet is one address.
The wallet table has \NWalletRoleRows{} wallet-role rows from the kept markets.
The kept markets hold \NTradesKeptMarkets{} trades.
We keep duplicate rows because no row identifier separates repeat fills from copies.

\section{Methods}\label{sec:methods}
\subsection{Evaluation price and pre-registration}\label{sec:methods-prereg}
% src: S-PAPER-2, S-PAPER-3
Each market has a scheduled close time $T$.
We set $T$ to the \texttt{close\_at} field.
Trades can continue after the close, often once the answer is known.
A price taken after the close can then contain the answer.
Then calibration looks better than it is.
In a one-month pilot of \PilotNMarkets{} markets, \PilotShareMarketsAfterClose{} had a trade after \texttt{close\_at}.
Those late trades are \PilotShareVolumeAfterClose{} of the pilot volume.
So we use the scheduled close and not the last trade.

\Cref{fig:timeline} shows where the evaluation prices sit.
The evaluation price \pOne{} is a mean price.
It averages the trades at the last timestamp at or before $T-\DaysShort$ day.
We drop a market if this price is older than \MaxStaleHours{} hours.
We also drop a market with fewer than \MinTrades{} trades before that time.
A price from a nearly empty market is not a crowd view.
We report Q1 with and without this filter.
The \DaysLong-day price \pSeven{} follows the same rule at $T-\DaysLong$ days.
It is used only in Q1b.

\begin{figure}[H]
\centering
\begin{tikzpicture}[font=\small]
\draw[thick,-{Stealth}] (0,0) -- (12.2,0);
\draw[thick] (0.2,-0.12) -- (0.2,0.12) node[above] {market opens};
\draw[thick,blue!60!black] (4.6,-0.2) -- (4.6,0.2) node[above] {$T-\DaysLong$ days};
\draw[thick,blue!60!black] (7.0,-0.2) -- (7.0,0.2) node[above] {$T-\DaysShort$ day};
\draw[thick] (9.2,-0.12) -- (9.2,0.12) node[above] {$T$ = \texttt{close\_at}};
\foreach \x in {10.0,10.5,11.0,11.5} \draw[red!70!black,thick] (\x,-0.15) -- (\x,0.15);
\node[below=6pt,blue!60!black,align=center] at (4.6,0) {price \pSeven\\(Q1b)};
\node[below=6pt,blue!60!black,align=center] at (7.0,0) {evaluation\\price \pOne};
\node[below=6pt,red!70!black,align=center] at (10.75,0) {trades after the close};
\end{tikzpicture}
\caption{Where the evaluation prices sit in a market's life (illustrative). Why: a price taken after the close can already contain the answer. Source: pre-registration, \texttt{analysis/prereg\_mvp1.md}.}
\label{fig:timeline}
\end{figure}

The criteria were fixed before any Q1 result.
\Cref{tab:prereg} lists them.
The verdict applies to the full-data run with the same parameters.

\begin{table}[H]
\centering
\caption{Pre-registered choices and Q1 criteria. Why: fixing the rules first keeps the result from choosing the rule. Source: \texttt{analysis/prereg\_mvp1.md}.}
\label{tab:prereg}
\footnotesize
\begin{tabularx}{\linewidth}{@{}lX@{}}
\toprule
Item & Value \\
\midrule
Close time $T$ & \texttt{close\_at} \\
\midrule
Evaluation price & mean price at the last timestamp (after the amendment below) at or before $T-\DaysShort$ day, at most \MaxStaleHours{} hours old \\
\midrule
Minimum trades & \MinTrades{} before $T-\DaysShort$ day. Q1 is also reported without it \\
\midrule
Bins and level & \NBins{} equal-width bins, $\alpha=\Alpha$ \\
\midrule
Primary metric & chi-square goodness of fit over the \NBins{} bins. The line is $p<\Alpha$ \\
\midrule
Per-bin call & binomial test per bin, Benjamini-Hochberg (BH) corrected at \Alpha{} \\
\midrule
Effect size & observed rate minus mean price per bin with a \CiLevel{} Wilson interval, and the Brier score \\
\midrule
Failure condition & no bin off after BH and chi-square $p\ge\Alpha$: no evidence of miscalibration at this $n$ \\
\bottomrule
\end{tabularx}
\end{table}

% src: S-PAPER-3
One rule changed after we saw a full-data Q1 result.
The original rule was the last price at or before $T-\DaysShort$ day.
Many trades share one block timestamp, so this rule left their order undefined.
The new rule is the mean price of the trades at the last timestamp.
We made the change for determinism, not because of the result.
It came after we saw a full-data Q1 result, so we report it as an amendment.
The amendment in \texttt{analysis/prereg\_mvp1.md} records that the verdict did not change.
All numbers in this paper come from the runs with the amended rule.

% src: S-PAPER-4
The MVP2 criteria were fixed before any MVP2 result.
\Cref{tab:prereg2} lists them.
Every sub-question uses the kept markets with the minimum-trade filter.
Each primary test is judged at \Alpha{} with no adjustment across the three.
A failed test reads ``no evidence \dots{} at this $n$''.

\begin{table}[H]
\centering
\caption{Pre-registered MVP2 criteria. Why: fixing the tests first keeps the result from choosing the test. Source: \texttt{analysis/prereg\_mvp2.md}.}
\label{tab:prereg2}
\footnotesize
\begin{tabularx}{\linewidth}{@{}lXX@{}}
\toprule
Question & Primary test and line & Effect size next to $p$ \\
\midrule
Q1a topic & chi-square test of independence, category by \{favorite won, favorite lost\}. Line: $p<\Alpha$ & Cram\'er's $V$ \\
\midrule
Q1b deadline & one-sample $t$-test of the per-market Brier difference ($\DaysLong$ days minus $\DaysShort$ day). Line: $p<\Alpha$ and mean difference above zero & mean difference with \CiLevel{} CI, Cohen's $d$ \\
\midrule
Q1c volume & Spearman $\rho$ of log volume against per-market Brier error. Line: $p<\Alpha$ and $\rho<0$ & $\rho$ with \CiLevel{} CI \\
\bottomrule
\end{tabularx}
\end{table}

One MVP2 choice was renamed after the review and before any code-version-\CodeVersion{} result.
Categories with too few markets merge into a group.
We first called that group ``Other''.
``Other'' is also a real category value.
So we renamed the group ``Merged''.
No criterion changed.

\subsection{Q1: calibration}\label{sec:methods-q1}
% src: S-PAPER-2, S-PAPER-3
We treat each outcome as a Bernoulli trial with success probability equal to the price.
So the Yes count in a bin is approximately binomial.
We form \NBins{} equal-width price bins, and we use equal-frequency bins as a check.
In bin $b$, $n_b$ is the market count, $\bar p_b$ the mean price, and $O_b$ the Yes count.
The null hypothesis is that the Yes rate equals the mean price.
The expected Yes count is $E_b=n_b\bar p_b$.
We test each bin with a binomial test.
We give a \CiLevel{} Wilson interval for the observed rate \citep{wilson1927probable}.
We correct the \NBins{} tests with Bonferroni and with the Benjamini-Hochberg (BH) procedure \citep{benjamini1995controlling}.
The overall test is chi-square goodness of fit on the bin by \{Yes, No\} table:
\begin{equation}
\chi^2=\sum_{b}\left[\frac{(O_b-E_b)^2}{E_b}+\frac{(O_b-E_b)^2}{n_b-E_b}\right],
\qquad w=\sqrt{\chi^2/n}.
\label{eq:chisq}
\end{equation}
Here $w$ is Cohen's effect size \citep{cohen1988statistical} and $n$ is the total market count.
We also report the Brier score $\frac1n\sum(p-y)^2$ \citep{brier1950verification}.
We add the log loss and the mean entropy $H(p)$ in bits.
We treat the market as a classifier that says Yes when $p>\FavoriteLine$.
We compare the price distributions of Yes- and No-resolved markets with a Kolmogorov-Smirnov (KS) test.
The KS statistic $D$ is the effect size.
For the longshot check we look at the two cheapest bins.

\subsection{Q1a, Q1b, Q1c: topic, horizon, volume}\label{sec:methods-q1abc}
% src: S-PAPER-4
For Q1a we use the field \texttt{category\_refined}.
A missing category is ``Unknown''.
A category with fewer than \MinCategoryN{} kept markets merges into ``Merged''.
The favorite is the side with price at least \FavoriteLine.
The favorite won if $\pOne\ge\FavoriteLine$ and $y=1$, or if $\pOne<\FavoriteLine$ and $y=0$.
The primary test is chi-square independence on category by \{favorite won, favorite lost\}.
Its effect size is Cram\'er's $V$ \citep{cramer1946mathematical}.
The secondary results are bin tests inside each category (bins with at least \MinBinN{} markets).
We correct them with Bonferroni and BH over all category-by-bin tests.
We also run an ANOVA on the Brier score across categories with its eta squared.
We add a permutation test with \QOneANPerm{} permutations and seed \PermSeed.

% src: S-PAPER-4
For Q1b the error of a market is its Brier score $(p-y)^2$.
We use the markets that opened at least \DaysLong{} days before $T$ and have a fresh \pSeven.
Each market is scored at both horizons.
The primary test is a one-sample $t$-test.
It uses the per-market difference, Brier at \DaysLong{} days minus Brier at \DaysShort{} day.
A positive mean means the market is better near the deadline.
We report the mean difference with its \CiLevel{} interval and Cohen's $d$.
Secondary results are the KS statistic between Yes and No prices at each horizon.
We add the difference in absolute error and the mean entropy $H(p)$.
The entropy falls by construction, because prices move toward zero or one near the deadline.
So the error change is the real test.

% src: S-PAPER-4
For Q1c the primary test is the Spearman correlation between $\log(1+\text{volume})$ and the per-market Brier error.
We give a \CiLevel{} interval from the Fisher transform.
The secondary result is the same correlation for the log trade count.
Volume and trade count cover all trades of the market in scope.

\subsection{Q2: wallet concentration}\label{sec:methods-q2}
% src: S-PAPER-2
We work with wallet volume in USDC, by role.
We report the share of role volume held by the top \TopOneLabel, \TopFiveLabel, and \TopTenLabel{} of wallets.
We draw the Lorenz curve and compute the Gini coefficient.
We compare our Gini with the dataset paper's as a check on our pipeline.
We test the log volume of makers against takers with a Welch t-test and a KS test.
We report Spearman's $\rho$ between trade count and volume.

% src: S-PAPER-2, S-PAPER-6
As a side analysis, we run k-means on wallet features.
We take the logarithm of the heavy-tailed features, then z-score them.
We choose $K$ with the elbow method.
We check stability by repeating the clustering with robust scaling.
We compare the two with the adjusted Rand index (ARI) \citep{hubert1985comparing}.
This analysis is exploratory.

\subsection{Q3: distributions}\label{sec:methods-q3}
% src: S-PAPER-2, S-PAPER-5
We summarise price, trade size, and wallet volume with their mean and median.
We plot histograms of the price \pOne{} and of trade size.
We fit a normal, an exponential, and a power law to trade size and wallet volume.
The power law is fitted on the top \TopTenLabel{} of values only, above a lower bound $x_{\min}$.
We measure each fit with the KS statistic $D$ against the data.
Trade sizes come from a deterministic sample of \NTradeSizeSample{} trades from the kept markets.

\subsection{Reproducibility}\label{sec:methods-repro}
% src: S-PAPER-5, S-PAPER-6
All data queries ran on the UTRGV Cradle cluster (SLURM, CPU only).
We used Python and DuckDB \citep{raasveldt2019duckdb}.
The MVP runs are \RunId{} (job \SlurmJobId), with code version \CodeVersion.
The hash \CodeSha{} is that of the Q1 script.
The hash of the whole script set is \CodeSetSha.
The per-script hashes are in \texttt{\CodeShaFile}.
The robustness run is \RobustRun, with code version \RobustCodeVersion{} (hash \RobustCodeSha).
An earlier review found three sources of small differences between reruns.
They were tie order inside one block, a random sample, and row order.
Code version \CodeVersion{} fixes these.
The pre-registered criteria did not change, and no verdict flipped.

\section{Results}\label{sec:results}
\subsection{Q1: is the crowd right?}\label{sec:res-q1}
% src: S-PAPER-3, S-PAPER-5
The pre-registered rule gives this verdict: \VerdictQOne{}
The chi-square statistic is \QOneChiSq{} with \QOneDf{} degrees of freedom.
The p-value is $p=\QOneP$, and the effect size is Cohen's $w=\QOneW$.
\Cref{fig:calib} shows the observed Yes rate against the mean price in each bin.
\Cref{tab:bins} gives the numbers.

\begin{table}[H]
\centering
\caption{Calibration by price bin, with the gap as observed rate minus mean price in percent points. Why: the gap and its interval show where prices are off. Source: run \RunId.}
\label{tab:bins}
\small
\setlength{\tabcolsep}{4pt}
\begin{tabular}{lrrrcrr}
\toprule
Bin & $n$ & Mean price & Observed & \CiLevel{} CI & Gap & BH $p$ \\
\midrule
\binrow{0--10\%}{ZeroTen}\\ \midrule
\binrow{10--20\%}{TenTwenty}\\ \midrule
\binrow{20--30\%}{TwentyThirty}\\ \midrule
\binrow{30--40\%}{ThirtyForty}\\ \midrule
\binrow{40--50\%}{FortyFifty}\\ \midrule
\binrow{50--60\%}{FiftySixty}\\ \midrule
\binrow{60--70\%}{SixtySeventy}\\ \midrule
\binrow{70--80\%}{SeventyEighty}\\ \midrule
\binrow{80--90\%}{EightyNinety}\\ \midrule
\binrow{90--100\%}{NinetyHundred}\\
\bottomrule
\end{tabular}
\end{table}

The observed rate is below the mean price in every bin except the top one.
After BH correction, \QOneBinsBh{} of the \NBins{} bins are off ($p<\Alpha$).
After Bonferroni correction, \QOneBinsBonferroni{} bins are off.
For example, in the bin with mean price \BinSeventyEightyPrice{}, the observed Yes rate is \BinSeventyEightyObs.
The gap is \BinSeventyEightyDiff{} (BH $p=\BinSeventyEightyBhP$).

The two cheapest bins are the longshot check.
In the cheapest bin the gap is \BinZeroTenDiff{} (BH $p=\BinZeroTenBhP$).
In the next bin the gap is \BinTenTwentyDiff{} (BH $p=\BinTenTwentyBhP$).
Cheap bets pay off less often than their price says.
The direction matches the horse-racing result \citep{snowberg2010favorite}.
Our gaps are per market and the dataset paper's are per trade.
So we do not compare their sizes.


\begin{figure}[H]
\centering
\begin{minipage}[b]{0.46\linewidth}\centering
\figorbox{figures/fig4_calibration}{\linewidth}
\end{minipage}\hfill
\begin{minipage}[b]{0.50\linewidth}\centering
\figorbox{figures/fig5_ecdf_yes_no}{\linewidth}
\end{minipage}
\caption{Observed Yes rate against mean price per bin with \CiLevel{} Wilson intervals (left), and ECDFs of \pOne{} by outcome (right). Why: points on the diagonal mean the price matches the outcome rate, and the farther apart the ECDFs, the better the price separates the outcomes. Source: run \RunId.}
\label{fig:calib}
\end{figure}

Next come the score measures.
The Brier score is \QOneBrier, the log loss is \QOneLogLoss, and the mean entropy is \QOneEntropy{} bits.
The Yes rate among kept markets is \QOneYesRate, and the mean price is \PriceOneDayMean.
At the \FavoriteLine{} line, the classifier accuracy is \QOneAccuracy.
Its precision is \QOnePrecision{} and its recall is \QOneRecall.
The KS statistic between Yes and No prices is $D=\QOneKsD$.
Without the minimum-trade filter there are \QOneNoMinN{} markets.
The Brier score is then \QOneNoMinBrier{}, with $\chi^2=\QOneNoMinChiSq$, $p=\QOneNoMinP$, and $w=\QOneNoMinW$.
The chi-square test also rejects without the filter.
With \QOneEqFreqBins{} equal-frequency bins, BH marks \QOneEqFreqBinsBh{} bins as off.

\subsection{Q1a: does accuracy depend on the topic?}\label{sec:res-q1a}
% src: S-PAPER-4, S-PAPER-5
The pre-registered rule gives this verdict: \VerdictQOneA{}
The chi-square statistic is \QOneAChiSq{} with \QOneADf{} degrees of freedom, and $p\,\QOneAP$ is a float-underflow bound.
The effect size is Cram\'er's $V=\QOneAV$.
The smallest expected count is \QOneAMinExpected.
\Cref{tab:cats} gives the favorite win rate and the mean Brier score in each of the \QOneANCategories{} groups.
Sports favorites win \CatSportsFavWin{} of the time and Crypto favorites win \CatCryptoFavWin.
The ANOVA on the Brier score gives $F=\QOneAAnovaF$ with $p\,\QOneAAnovaP$ and $\eta^2=\QOneAEtaSq$.
The permutation $p$ is at the floor, $1/(\QOneANPerm+1)$, for \QOneANPerm{} permutations.

\begin{table}[H]
\centering
\caption{By topic: market count, favorite win rate, mean Brier score, and the longshot gap (observed rate minus mean price, in percent points) in the two cheapest bins. Why: a topic that differs from the all-topics row shows where the price misleads. BH covers all topic-by-bin tests, and the bins for all topics. Source: run \RunId.}
\label{tab:cats}
\footnotesize
\setlength{\tabcolsep}{3pt}
\begin{tabular}{lrrrrrrrrr}
\toprule
 & & & & \multicolumn{3}{c}{Cheapest bin} & \multicolumn{3}{c}{Next bin} \\
\cmidrule(lr){5-7}\cmidrule(lr){8-10}
Topic & Markets & Fav. win & Brier & $n$ & Gap & BH $p$ & $n$ & Gap & BH $p$ \\
\midrule
\topicrow{Sports}{Sports}{CatSportsBinZeroTen}{CatSportsBinTenTwenty}\\ \midrule
\topicrow{Crypto}{Crypto}{CatCryptoBinZeroTen}{CatCryptoBinTenTwenty}\\ \midrule
\topicrow{Politics}{Politics}{CatPoliticsBinZeroTen}{CatPoliticsBinTenTwenty}\\ \midrule
\topicrow{Finance}{Finance}{CatFinanceBinZeroTen}{CatFinanceBinTenTwenty}\\ \midrule
\topicrow{Other}{Other}{CatOtherBinZeroTen}{CatOtherBinTenTwenty}\\ \midrule
\topicrow{Culture}{Culture}{CatCultureBinZeroTen}{CatCultureBinTenTwenty}\\ \midrule
\topicrow{Sci-Tech}{SciTech}{CatSciTechBinZeroTen}{CatSciTechBinTenTwenty}\\ \midrule
\topicrow{Price Action}{PriceAction}{CatPriceActionBinZeroTen}{CatPriceActionBinTenTwenty}\\ \midrule
All topics & \QOneAN & \CatAllFavWin & \QOneBrier & \BinZeroTenN & \BinZeroTenDiff & \BinZeroTenBhP & \BinTenTwentyN & \BinTenTwentyDiff & \BinTenTwentyBhP \\
\bottomrule
\end{tabular}
\end{table}

\begin{figure}[H]
\centering
\figorbox{figures/fig9_calibration_by_category}{\linewidth}
\caption{Observed Yes rate against mean price per bin, one panel per topic, on shared axes. Why: a panel off the diagonal shows a topic where the price is miscalibrated. Source: run \RunId.}
\label{fig:bycat}
\end{figure}

% src: S-PAPER-7
\Cref{fig:bycat} shows the calibration by topic.
The bin tests inside topics cover \QOneABinsTested{} topic-by-bin cells.
BH marks \QOneABinsBh{} of them as off, and Bonferroni marks \QOneABinsBonferroni.
\Cref{tab:cats} also shows the two cheapest bins by topic.
In the next bin, cheap bets in Crypto and Politics pay off less often than priced.
The gaps are \CatCryptoBinTenTwentyDiff{} and \CatPoliticsBinTenTwentyDiff{}, with BH $p=\CatCryptoBinTenTwentyBhP$ and $p=\CatPoliticsBinTenTwentyBhP$.
In Sports the gap is \CatSportsBinTenTwentyDiff{} (BH $p=\CatSportsBinTenTwentyBhP$).
In the cheapest bin, only Crypto is off after BH (BH $p=\CatCryptoBinZeroTenBhP$).
The direction for Crypto and Politics matches the horse-racing result \citep{snowberg2010favorite}.

\subsection{Q1b: does the market get more accurate near the deadline?}\label{sec:res-q1b}
% src: S-PAPER-4, S-PAPER-5
The pre-registered rule gives this verdict: \VerdictQOneB{}
The test uses \QOneBN{} markets.
The mean Brier difference is \QOneBMeanDiff{} with a \CiLevel{} interval of [\QOneBCiLow, \QOneBCiHigh].
The statistic is $t=\QOneBT$ with \QOneBDf{} degrees of freedom and $p=\QOneBP$.
The effect size is Cohen's $d=\QOneBD$.
The absolute error falls by \QOneBAbsMeanDiff{} on average (Cohen's $d=\QOneBAbsD$).
\Cref{fig:horizon} compares the two horizons.
The mean Brier score is \QOneBBrierSevenDays{} at \DaysLong{} days and \QOneBBrierOneDay{} at \DaysShort{} day.
The KS statistic between Yes and No prices is $D=\QOneBKsDSevenDays$ at \DaysLong{} days and $D=\QOneBKsDOneDay$ at \DaysShort{} day.
The mean entropy falls by \QOneBEntropyDiff{} bits.
This fall is expected by construction, so we do not count it as evidence.

\begin{figure}[H]
\centering
\figorbox{figures/fig10_ecdf_7d_vs_1d}{0.9\linewidth}
\caption{ECDFs of the price for Yes- and No-resolved markets at \DaysLong{} days and at \DaysShort{} day before $T$. Why: the farther apart the two curves, the better the price separates the outcomes. Source: run \RunId.}
\label{fig:horizon}
\end{figure}

\subsection{Q1c: are busier markets more accurate?}\label{sec:res-q1c}
% src: S-PAPER-4, S-PAPER-5
The pre-registered rule gives this verdict: \VerdictQOneC{}
The correlation between log volume and Brier error is $\rho=\QOneCVolumeRho$
(\CiLevel{} CI [\QOneCVolumeCiLow, \QOneCVolumeCiHigh], $p\,\QOneCVolumeP$).
The sign is positive, so the pre-registered direction did not appear.
The log trade count gives $\rho=\QOneCTradesRho$ (\CiLevel{} CI [\QOneCTradesCiLow, \QOneCTradesCiHigh], $p\,\QOneCTradesP$).
\Cref{fig:volume} shows the relation.

\begin{figure}[H]
\centering
\figorbox{figures/fig11_volume_vs_error}{0.9\linewidth}
\caption{Per-market Brier error against volume and trade count on log axes, with the mean error per decile. Why: a rising line means busier markets have larger errors. Source: run \RunId.}
\label{fig:volume}
\end{figure}

\subsection{Q2: who is the crowd?}\label{sec:res-q2}
% src: S-PAPER-5
In both roles, the top \TopOneLabel{} of wallets hold most of the volume.
\Cref{tab:conc} gives the shares and the Gini coefficient by role.
\Cref{fig:lorenz} shows the Lorenz curves.
We describe these numbers without a verdict.

\begin{table}[H]
\centering
\caption{Volume concentration by role, as the share of role volume held by the top wallets. Why: the top shares show how few wallets hold most of the volume. Source: run \RunId.}
\label{tab:conc}
\small
\begin{tabular}{lrrrrrrr}
\toprule
Role & Wallets & Mean (USDC) & Median (USDC) & Top \TopOneLabel{} & Top \TopFiveLabel{} & Top \TopTenLabel{} & Gini \\
\midrule
Maker & \MakerNWallets & \MakerMeanVolume & \MakerMedianVolume & \MakerTopOneShare & \MakerTopFiveShare & \MakerTopTenShare & \MakerGini \\
\midrule
Taker & \TakerNWallets & \TakerMeanVolume & \TakerMedianVolume & \TakerTopOneShare & \TakerTopFiveShare & \TakerTopTenShare & \TakerGini \\
\bottomrule
\end{tabular}
\end{table}

\begin{figure}[H]
\centering
\figorbox{figures/fig7_lorenz}{0.5\linewidth}
\caption{Lorenz curves of wallet volume for makers and takers. Why: a curve far below the equality line means few wallets hold most of the volume. Source: run \RunId.}
\label{fig:lorenz}
\end{figure}

The mean volume is far above the median volume in both roles.
The dataset paper reports a Gini of \PaperGiniMaker{} for makers and \PaperGiniTaker{} for takers \citep{qin2026polymarketv1}.
Ours are \MakerGini{} and \TakerGini.
We use the paper's values as a reference only.
The paper covers the full data, and our wallets come from the kept markets.

The Welch test on log volume gives $t=\RoleWelchT$ with $p\,\RoleWelchP$, a float-underflow bound.
The effect size is Cohen's $d=\RoleCohenD$, and the KS statistic is $D=\RoleKsD$.
Trade count and volume have a Spearman correlation of $\rho=\SpearmanTradesVolumeRho$.

In the side analysis, the elbow method picks $K=\KmeansK$.
The sample has \KmeansNSample{} wallet-role rows, and a wallet in both roles gives two rows.
The \KmeansK{} clusters have \ClusterZeroN, \ClusterOneN, \ClusterTwoN, and \ClusterThreeN{} rows.
Their median volumes run from \ClusterTwoMedianVolume{} to \ClusterOneMedianVolume{} USDC.
Clusters from z-score and robust scaling agree (ARI $=\KmeansAri$).
Cramér's $V$ \citep{cramer1946mathematical} between cluster and market category is \KmeansCategoryV.
An ANOVA of mean trade price across clusters gives $\eta^2=\KmeansPriceEtaSq$.
Its permutation $p$ is at the floor, $1/(\KmeansNPerm+1)$.

\subsection{Q3: what does the data look like?}\label{sec:res-q3}
% src: S-PAPER-5
We describe the distributions without a verdict.
\Cref{fig:hist} shows the price histogram and the trade-size histogram.

\begin{figure}[H]
\centering
\figorbox{figures/fig1_fig2_hist}{\linewidth}
\caption{Histogram of the price \pOne{} over kept markets (left) and of trade size on a log axis (right). Why: the price shape shows where the crowd is confident, and the trade-size shape shows how uneven bets are. Source: run \RunId.}
\label{fig:hist}
\end{figure}

The price \pOne{} has a mean of \PriceOneDayMean{} and a median of \PriceOneDayMedian.
Prices pile up near both ends of the scale.
The trade size has a mean of \TradeSizeMean{} USDC and a median of \TradeSizeMedian{} USDC.
So trade sizes are strongly right skewed.

\Cref{tab:fits} gives the fit statistics.
The KS statistic $D$ is the effect size.
KS p-values are optimistic here because the parameters are fitted on the same data.
So we do not rely on them.
The normal fit is poor for both variables.
The power-law fit covers the tail only, so its $D$ is not directly comparable.
The exponential fit has $D=\TradeSizeExponentialKsD$ for trade size and $D=\WalletVolumeExponentialKsD$ for wallet volume.

\begin{table}[H]
\centering
\caption{Distribution fits, with the power law fitted on the tail from $x_{\min}$. Why: a small $D$ means the fitted curve is close to the data. Source: run \RunId.}
\label{tab:fits}
\small
\begin{tabular}{lrrrr}
\toprule
Variable & Power-law $\alpha$ & $x_{\min}$ (USDC) & KS $D$, power law & KS $D$, normal \\
\midrule
Trade size & \TradeSizePowerAlpha & \TradeSizePowerXmin & \TradeSizePowerKsD & \TradeSizeNormalKsD \\
\midrule
Wallet volume & \WalletVolumePowerAlpha & \WalletVolumePowerXmin & \WalletVolumePowerKsD & \WalletVolumeNormalKsD \\
\bottomrule
\end{tabular}
\end{table}

\subsection{Robustness}\label{sec:res-robust}
% src: S-PAPER-9, S-PAPER-8
Two checks look at limits of the Q1 setup.
Both are descriptive.
They add no primary test and no verdict.
The pre-registered Q1 verdict stays as it is.
The sensitivity row below sits next to the original Q1 row and does not replace it.
Both checks use the \QOneN{} kept markets.

% src: S-PAPER-9
Check R1 asks whether the Q1 result holds when markets without a clean payout are removed.
The resolution file gives each market a payout vector.
A market is clean when one outcome pays one and the other pays zero.
Of the kept markets, \RobROneCleanN{} are clean (\RobROneCleanShare).
\RobROneSplitN{} have a split payout (\RobROneSplitShare).
\RobROneOtherN{} are of another kind.
\RobROneMissingN{} is missing from the file.
Our label marks Yes on \RobROneSplitYesOne{} of the split markets and No on \RobROneSplitYesZero.
Our winner label agrees with the resolution file on \RobROneAgreement{} of the \RobROneNClean{} clean markets.
The \CiLevel{} Wilson CI is [\RobROneAgreementLow, \RobROneAgreementHigh].
\RobROneNMismatch{} clean markets disagree.
\Cref{fig:r1} shows the payout classes and the calibration points with and without the unclear markets.
\Cref{tab:r1} puts the sensitivity row next to the original row.
We remove split, other, and mismatching markets, which is \RobROneNExcluded{} markets.
The Q1 rule result is the same in both rows: \RobROneRuleOriginal{}
The effect size changes by \RobROneDeltaW.

\begin{figure}[H]
\centering
\figorbox{figures/fig_r1_split_payout}{0.9\linewidth}
\caption{Kept markets by payout class (log scale), and the Q1 calibration points with and without the unclear markets. Why: points that overlap mean the unclear markets do not move the result. Source: run \RobustRun.}
\label{fig:r1}
\end{figure}

\begin{table}[H]
\centering
\caption{Q1 original row and R1 sensitivity row. Why: equal rows mean the unclear markets do not drive the Q1 result. Source: runs \RunId{} and \RobustRun.}
\label{tab:r1}
\small
\begin{tabular}{lrrrrrr}
\toprule
Row & Markets & $\chi^2$ & $p$ & Cohen's $w$ & Brier & Bins off (BH) \\
\midrule
Q1 original & \QOneN & \QOneChiSq & \QOneP & \QOneW & \QOneBrier & \QOneBinsBh \\
\midrule
Sensitivity & \RobROneSensN & \RobROneSensChiSq & \RobROneSensP & \RobROneSensW & \RobROneSensBrier & \RobROneSensBinsBh \\
\bottomrule
\end{tabular}
\end{table}

% src: S-PAPER-9
Check R2 asks how far the price moves from $T-\DaysShort$ day to the last trade before $T$.
Call the mean price at the last timestamp before $T$ \pT.
The median absolute change from \pOne{} to \pT{} is \RobRTwoMedianAbsDelta.
The change is at least \DeltaA{} in \RobRTwoShareDeltaAtLeastTen{} of markets.
It is at least \DeltaB{} in \RobRTwoShareDeltaAtLeastTwentyFive{} and at least \DeltaC{} in \RobRTwoShareDeltaAtLeastFifty.
The price crosses \FavoriteLine{} in \RobRTwoShareCrossing{} of markets.
This compares the two end points and not the path.
\RobRTwoShareNoTrade{} of markets have no trade in the last day, and for them \pT{} equals \pOne.
\Cref{fig:r2} shows the distribution of the absolute change.
\Cref{tab:r2cat} gives the change and the Brier difference by topic.
There is no p-value in this table.

\begin{table}[H]
\centering
\caption{R2 by topic: market count, median absolute change from \pOne{} to \pT, and mean Brier difference (\pT{} minus \pOne). Why: topics with a large change show where \pOne{} is least like the close price. Source: run \RobustRun.}
\label{tab:r2cat}
\small
\begin{tabular}{lrrr}
\toprule
Topic & Markets & Median $|\Delta p|$ & Mean Brier difference \\
\midrule
Sports & \RobRTwoCatSportsN & \RobRTwoCatSportsMedianAbsDelta & \RobRTwoCatSportsMeanDiff \\ \midrule
Crypto & \RobRTwoCatCryptoN & \RobRTwoCatCryptoMedianAbsDelta & \RobRTwoCatCryptoMeanDiff \\ \midrule
Politics & \RobRTwoCatPoliticsN & \RobRTwoCatPoliticsMedianAbsDelta & \RobRTwoCatPoliticsMeanDiff \\ \midrule
Finance & \RobRTwoCatFinanceN & \RobRTwoCatFinanceMedianAbsDelta & \RobRTwoCatFinanceMeanDiff \\ \midrule
Other & \RobRTwoCatOtherN & \RobRTwoCatOtherMedianAbsDelta & \RobRTwoCatOtherMeanDiff \\ \midrule
Culture & \RobRTwoCatCultureN & \RobRTwoCatCultureMedianAbsDelta & \RobRTwoCatCultureMeanDiff \\ \midrule
Sci-Tech & \RobRTwoCatSciTechN & \RobRTwoCatSciTechMedianAbsDelta & \RobRTwoCatSciTechMeanDiff \\ \midrule
Price Action & \RobRTwoCatPriceActionN & \RobRTwoCatPriceActionMedianAbsDelta & \RobRTwoCatPriceActionMeanDiff \\
\bottomrule
\end{tabular}
\end{table}

\begin{figure}[H]
\centering
\figorbox{figures/fig_r2_price_path}{0.7\linewidth}
\caption{ECDF of the absolute price change from \pOne{} to \pT. Why: a curve near one at small values means \pOne{} is close to the last price. Source: run \RobustRun.}
\label{fig:r2}
\end{figure}

The price \pT{} is the price just before the close.
So its Brier score is not a one-day-ahead forecast.
The mean Brier score is \RobRTwoAllBrierOneDay{} at \pOne{} and \RobRTwoAllBrierAtClose{} at \pT.
The paired difference (\pT{} minus \pOne) is \RobRTwoAllMeanDiff{}
(\CiLevel{} CI [\RobRTwoAllCiLow, \RobRTwoAllCiHigh], $p\,\RobRTwoAllP$, Cohen's $d_z=\RobRTwoAllDz$).
A negative difference means \pT{} is closer to the outcome than \pOne.
On the \RobRTwoTradedN{} markets with a trade in the last day, the difference is \RobRTwoTradedMeanDiff.
The \CiLevel{} CI is [\RobRTwoTradedCiLow, \RobRTwoTradedCiHigh], and $d_z=\RobRTwoTradedDz$.

\section{Discussion}\label{sec:discussion}
% src: S-PAPER-5, S-PAPER-8
The answer to Q1 is no.
At this $n$, the price is not a calibrated probability as a whole.
The observed Yes rate sits below the price in all but the top bin.
The effect size is Cohen's $w=\QOneW$, and the largest bin gap is \BinSeventyEightyDiff{} (\Cref{tab:bins}).
The longshot gap points the same way as the horse-racing bias.
Accuracy depends on the topic.
Favorites win \CatSportsFavWin{} of the time in Sports and \CatCryptoFavWin{} in Crypto.
The longshot gap shows in Crypto and Politics.
It is not detected in Sports (gap \CatSportsBinTenTwentyDiff, BH $p=\CatSportsBinTenTwentyBhP$).
The market is closer to the outcome \DaysShort{} day before the close than \DaysLong{} days before.
We find no evidence that busier markets are more accurate at this $n$.
The correlation is positive, the opposite of the pre-registered direction.
\todo{owner: interpretation of why prices run high, why Sports differs from Crypto, why busier markets show larger errors, and what the gaps mean in practice for a reader of the price}

The answer to Q2 is descriptive.
Volume is concentrated in few wallets, in both roles.
Our Gini values are \MakerGini{} and \TakerGini, against \PaperGiniMaker{} and \PaperGiniTaker{} in the dataset paper.
The populations differ.
We measure concentration of volume.
We do not test who sets the prices.
The k-means clusters are a side view and carry no claim.

The answer to Q3 is descriptive too.
Prices cluster near the ends, and trade sizes are strongly right skewed.
\todo{owner: wider interpretation or policy implications, if the course report needs them}

\section{Limitations}\label{sec:limits}
% src: S-PAPER-1, S-PAPER-2, S-PAPER-3, S-PAPER-6
Some markets have no usable winner label, and we drop them.
In our full run this is \NNoOutcome{} of \NMarketsInScope{} markets (\PctNoOutcome).
Our rule is wider than the dataset paper's.
It also drops markets whose Yes-side label never traded.
The paper says only a small subset lacks a winner.

Markets keep trading after \texttt{close\_at}.
We avoid this by pricing one day before the scheduled close.
But the choice rests on a one-month pilot.

Sports is the largest topic in scope and in the kept markets.
It is \SportsShareInScope{} of the \NMarketsInScope{} markets in scope and \SportsShareKept{} of the kept markets.
Its survival is \CatSportsSurvival{} against \PctKept{} overall.
So Sports is slightly under-represented against its scope.
Sports loses markets mainly to no price (\CatSportsDropNoPriceN) or a stale price (\CatSportsDropStalePriceN).
Both are measured \DaysShort{} day before the close.
It loses almost none to the Up/Down rule (\CatSportsDropUpdownN).
The dataset paper notes that most Sports markets are Neg-Risk or Up/Down \citep{qin2026polymarketv1}.
Neg-Risk markets are outside the layer we use.

Our selection is narrow.
Only \NMarketsKept{} of \NMarketsInScope{} markets pass the filters.
\NNoPrice{} markets have no price \DaysShort{} day before the close (\PctNoPrice).
The minimum-trade filter is a choice, and we show the result without it.
Markets on one event move together.
The binomial and chi-square tests assume independent markets, so the p-values are optimistic.
The large $n$ also makes p-values underflow to zero in places.
We therefore read the effect sizes and the intervals first.

The k-means result is exploratory.
The stability check gives ARI $=\KmeansAri$ between two scalings.
This does not make the clusters real.
A wallet is an address, and one owner may control many addresses.
Volume concentration is not the same as the power to set prices.

The Q1b result rests on a subset.
It uses \QOneBN{} of the \NMarketsKept{} kept markets.
These are open at least \DaysLong{} days and have a fresh \pSeven.
These markets are long and active.
In Q1c the volume and trade count cover all trades in scope, not only those before $T-\DaysShort$ day.
The Brier error is larger for contested markets by construction, so $\rho$ mixes activity with contest.
The three primary tests of MVP2 are each judged at \Alpha{}.
We make no adjustment across them.
``Other'' is a real topic.
So we renamed the merge group to ``Merged'' (\Cref{sec:methods-prereg}).

We changed one pre-registered rule after seeing a full Q1 result.
This is the tie rule in \Cref{sec:methods-prereg}.
We report it as an amendment.
The amendment records that the verdict did not change.

The two robustness checks in \Cref{sec:res-robust} speak to two limits.
In R1, only \RobROneSplitN{} kept markets have a split payout and \RobROneMissingN{} is missing from the file.
Our label agrees with the resolution file on \RobROneAgreement{} of clean markets.
The Q1 rule result is the same without the unclear markets.
In R2, the price moves by at least \DeltaA{} in \RobRTwoShareDeltaAtLeastTen{} of markets.
This is the move between \pOne{} and the last trade.
So \pOne{} is not the price at the close.
Also \RobRTwoShareNoTrade{} of markets have no trade in the last day.

\section{Conclusion}\label{sec:conclusion}
% src: S-PAPER-5, S-PAPER-8
We tested whether a Polymarket price can be read as a probability.
On \NMarketsKept{} resolved Yes/No markets, the answer is not exactly.
The prices are not calibrated as a whole.
Cheap bets pay off less often than priced, which matches the horse-racing direction.
Accuracy depends on the topic and improves toward the deadline.
We find no evidence that busier markets are more accurate.
The top wallets hold most of the volume, and trade sizes are strongly right skewed.
We describe the last two results without a verdict.
The payout check leaves the Q1 result in place.
\todo{owner: closing sentence on the value of the study, and future work}

\begingroup\raggedright
\bibliographystyle{plainnat}
\bibliography{references}
\endgroup

\end{document}
